\section*{Lecture 3: Uncertainty}
\addcontentsline{toc}{section}{Lecture 3: Uncertainty}
\clearpage
\SlideHeader{Lecture 3: Uncertainty}{001}{表紙}
\SlideImage{1}{../Materials/2026/3DM_03_Uncertainty_SS26.pdf}
\begin{adjustbox}{max totalsize={\textwidth}{0.675\textheight},center}
\begin{minipage}{\textwidth}\scriptsize
\begin{minipage}[t]{0.485\textwidth}
\NoteSection{日本語訳}
\begin{itemize}
\item Data Driven Decision Making
\item 第3回 - 不確実性
\end{itemize}
\end{minipage}\hfill
\begin{minipage}[t]{0.485\textwidth}
\NoteSection{注記}
{\small このページは表紙・目次・連絡先などの事務情報であるため、解説と確認問題は付けない。}
\end{minipage}
\end{minipage}
\end{adjustbox}

\clearpage
\SlideHeader{Lecture 3: Uncertainty}{002}{復習と本日}
\SlideImage{2}{../Materials/2026/3DM_03_Uncertainty_SS26.pdf}
\begin{adjustbox}{max totalsize={\textwidth}{0.675\textheight},center}
\begin{minipage}{\textwidth}\scriptsize
\begin{minipage}[t]{0.485\textwidth}
\NoteSection{日本語訳}
\begin{itemize}
\item Business Analytics の全プロセス。
\item 情報モデルと意思決定モデル。
\item データが最初に来る。
\item 正しい情報モデルを見つけることは難しい。
\item より精緻な情報モデルにより意思決定の複雑性が増す。
\item 情報モデルにおける不確実性。
\item 不確実性をどのようにモデル化できるか。
\item 不確実性の下でどのように意思決定できるか。
\item ケーススタディ: 不確実な旅行時間を持つ配送ルーティング。
\end{itemize}
\NoteSection{試験対策ポイント}
\begin{itemize}
\item 不確実性は、需要・リソース・環境の変化として整理できる。
\item 頻度、範囲、破壊力、予測可能性で性質が異なる。
\end{itemize}
\end{minipage}\hfill
\begin{minipage}[t]{0.485\textwidth}
\NoteSection{解説}
このページでは「復習と本日」を、講義全体の流れの中で位置づける。スライドの箇条書きは短いが、試験ではその背後にある概念、現実例、モデル上の意味を文章で説明できる必要がある。\par
不確実性とは、意思決定時点では将来の状態や入力値が完全には分からないことである。配送であれば、注文がどこから来るか、交通がどれだけ混むか、サービス時間が何分か、ドライバーが利用可能かどうかが事前には確定していない。\par
不確実性は一種類ではない。需要不確実性は、顧客からの注文の発生時刻、場所、量に関する不確実性である。リソース不確実性は、車両故障、ドライバーのキャンセル、バッテリー残量、積載能力に関する不確実性である。環境不確実性は、交通、天候、道路工事、事故、駐車可能性に関する不確実性である。\par
不確実性は、頻度、範囲、破壊力、予測可能性によって性質が異なる。例えば日々の交通変動は頻繁に起きるが、過去データからある程度予測できる。一方で、大規模事故や道路閉鎖は低頻度だが、発生すれば計画を大きく壊す。\par
不確実性を無視して平均値だけで計画すると、平均的には良く見えるが実行時に頻繁に遅れる計画になる可能性がある。特に複数顧客を連続して訪問するルーティングでは、一つの遅延が後続訪問に波及するため、単独の平均時間以上にリスクが大きい。\par
試験答案では、不確実性の例を挙げるだけでなく、それが情報モデルと意思決定モデルにどう影響するかを書くとよい。情報モデルでは確率分布や区間として表現し、意思決定モデルでは期待値、リスク、ロバスト性、再最適化などを考える。\par
\NoteSection{確認問題と解答}
\begin{enumerate}
\item \textbf{Slide 002「復習と本日」の中心概念を、専門用語を使わずに説明せよ。}\\このスライドでは、復習と本日 を通じて 不確実性は、需要・リソース・環境の変化として整理できる。 ことを理解する必要がある。試験答案では、まず概念の定義を一文で書き、次に講義の例へ対応付け、最後に意思決定モデルにどのような影響を与えるかを述べる。単語の列挙だけでは不十分であり、データがどのように情報モデルへ変換され、その情報モデルがどのように決定変数・目的関数・制約へ接続されるかまで説明する。
\item \textbf{Slide 002「復習と本日」を、配送またはTSPの具体例で説明せよ。}\\注文がいつ来るか、ドライバーが稼働するか、道路が混むかは意思決定時点では確定しない。これが不確実性である。 答案では、例を出した後に、それがどのモデル要素に対応するかを明示する。例えば、顧客集合や旅行時間行列は情報モデル、訪問順序や車両割当は意思決定モデル、実際の配送指示は実装である。
\end{enumerate}
\end{minipage}
\end{minipage}
\end{adjustbox}

\clearpage
\SlideHeader{Lecture 3: Uncertainty}{003}{アジェンダ}
\SlideImage{3}{../Materials/2026/3DM_03_Uncertainty_SS26.pdf}
\begin{adjustbox}{max totalsize={\textwidth}{0.675\textheight},center}
\begin{minipage}{\textwidth}\scriptsize
\begin{minipage}[t]{0.485\textwidth}
\NoteSection{日本語訳}
\begin{itemize}
\item 不確実性。
\item 不確実性のモデリング。
\item 不確実性下の意思決定。
\item ケーススタディ: 都市物流におけるロバスト車両ルーティング。
\end{itemize}
\end{minipage}\hfill
\begin{minipage}[t]{0.485\textwidth}
\NoteSection{注記}
{\small このページは表紙・目次・連絡先などの事務情報であるため、解説と確認問題は付けない。}
\end{minipage}
\end{minipage}
\end{adjustbox}

\clearpage
\SlideHeader{Lecture 3: Uncertainty}{004}{導入}
\SlideImage{4}{../Materials/2026/3DM_03_Uncertainty_SS26.pdf}
\begin{adjustbox}{max totalsize={\textwidth}{0.675\textheight},center}
\begin{minipage}{\textwidth}\scriptsize
\begin{minipage}[t]{0.485\textwidth}
\NoteSection{日本語訳}
\begin{itemize}
\item 例
\begin{itemize}
\item 株式市場
\item 需要
\item 供給
\item イノベーション
\item 政治
\end{itemize}
\item 違い
\begin{itemize}
\item 頻度
\item 範囲
\item 破壊力
\item 予測可能性
\end{itemize}
\end{itemize}
\NoteSection{試験対策ポイント}
\begin{itemize}
\item 不確実性は、需要・リソース・環境の変化として整理できる。
\item 頻度、範囲、破壊力、予測可能性で性質が異なる。
\end{itemize}
\end{minipage}\hfill
\begin{minipage}[t]{0.485\textwidth}
\NoteSection{解説}
このページでは「導入」を、講義全体の流れの中で位置づける。スライドの箇条書きは短いが、試験ではその背後にある概念、現実例、モデル上の意味を文章で説明できる必要がある。\par
不確実性とは、意思決定時点では将来の状態や入力値が完全には分からないことである。配送であれば、注文がどこから来るか、交通がどれだけ混むか、サービス時間が何分か、ドライバーが利用可能かどうかが事前には確定していない。\par
不確実性は一種類ではない。需要不確実性は、顧客からの注文の発生時刻、場所、量に関する不確実性である。リソース不確実性は、車両故障、ドライバーのキャンセル、バッテリー残量、積載能力に関する不確実性である。環境不確実性は、交通、天候、道路工事、事故、駐車可能性に関する不確実性である。\par
不確実性は、頻度、範囲、破壊力、予測可能性によって性質が異なる。例えば日々の交通変動は頻繁に起きるが、過去データからある程度予測できる。一方で、大規模事故や道路閉鎖は低頻度だが、発生すれば計画を大きく壊す。\par
不確実性を無視して平均値だけで計画すると、平均的には良く見えるが実行時に頻繁に遅れる計画になる可能性がある。特に複数顧客を連続して訪問するルーティングでは、一つの遅延が後続訪問に波及するため、単独の平均時間以上にリスクが大きい。\par
試験答案では、不確実性の例を挙げるだけでなく、それが情報モデルと意思決定モデルにどう影響するかを書くとよい。情報モデルでは確率分布や区間として表現し、意思決定モデルでは期待値、リスク、ロバスト性、再最適化などを考える。\par
\NoteSection{確認問題と解答}
\begin{enumerate}
\item \textbf{Slide 004「導入」の中心概念を、専門用語を使わずに説明せよ。}\\このスライドでは、導入 を通じて 不確実性は、需要・リソース・環境の変化として整理できる。 ことを理解する必要がある。試験答案では、まず概念の定義を一文で書き、次に講義の例へ対応付け、最後に意思決定モデルにどのような影響を与えるかを述べる。単語の列挙だけでは不十分であり、データがどのように情報モデルへ変換され、その情報モデルがどのように決定変数・目的関数・制約へ接続されるかまで説明する。
\item \textbf{Slide 004「導入」を、配送またはTSPの具体例で説明せよ。}\\注文がいつ来るか、ドライバーが稼働するか、道路が混むかは意思決定時点では確定しない。これが不確実性である。 答案では、例を出した後に、それがどのモデル要素に対応するかを明示する。例えば、顧客集合や旅行時間行列は情報モデル、訪問順序や車両割当は意思決定モデル、実際の配送指示は実装である。
\end{enumerate}
\end{minipage}
\end{minipage}
\end{adjustbox}

\clearpage
\SlideHeader{Lecture 3: Uncertainty}{005}{モビリティと物流における不確実性}
\SlideImage{5}{../Materials/2026/3DM_03_Uncertainty_SS26.pdf}
\begin{adjustbox}{max totalsize={\textwidth}{0.675\textheight},center}
\begin{minipage}{\textwidth}\scriptsize
\begin{minipage}[t]{0.485\textwidth}
\NoteSection{日本語訳}
\begin{itemize}
\item 顧客: 依頼、需要量、サービス時間、利用可能性、行動。
\item ドライバー: 利用可能性、行動、航続距離、スキル。
\item 環境: 旅行時間、駐車、ネットワーク利用可能性。
\end{itemize}
\NoteSection{試験対策ポイント}
\begin{itemize}
\item 不確実性は、需要・リソース・環境の変化として整理できる。
\item 頻度、範囲、破壊力、予測可能性で性質が異なる。
\end{itemize}
\end{minipage}\hfill
\begin{minipage}[t]{0.485\textwidth}
\NoteSection{解説}
このページでは「モビリティと物流における不確実性」を、講義全体の流れの中で位置づける。スライドの箇条書きは短いが、試験ではその背後にある概念、現実例、モデル上の意味を文章で説明できる必要がある。\par
不確実性とは、意思決定時点では将来の状態や入力値が完全には分からないことである。配送であれば、注文がどこから来るか、交通がどれだけ混むか、サービス時間が何分か、ドライバーが利用可能かどうかが事前には確定していない。\par
不確実性は一種類ではない。需要不確実性は、顧客からの注文の発生時刻、場所、量に関する不確実性である。リソース不確実性は、車両故障、ドライバーのキャンセル、バッテリー残量、積載能力に関する不確実性である。環境不確実性は、交通、天候、道路工事、事故、駐車可能性に関する不確実性である。\par
不確実性は、頻度、範囲、破壊力、予測可能性によって性質が異なる。例えば日々の交通変動は頻繁に起きるが、過去データからある程度予測できる。一方で、大規模事故や道路閉鎖は低頻度だが、発生すれば計画を大きく壊す。\par
不確実性を無視して平均値だけで計画すると、平均的には良く見えるが実行時に頻繁に遅れる計画になる可能性がある。特に複数顧客を連続して訪問するルーティングでは、一つの遅延が後続訪問に波及するため、単独の平均時間以上にリスクが大きい。\par
試験答案では、不確実性の例を挙げるだけでなく、それが情報モデルと意思決定モデルにどう影響するかを書くとよい。情報モデルでは確率分布や区間として表現し、意思決定モデルでは期待値、リスク、ロバスト性、再最適化などを考える。\par
\NoteSection{確認問題と解答}
\begin{enumerate}
\item \textbf{Slide 005「モビリティと物流における不確実性」の中心概念を、専門用語を使わずに説明せよ。}\\このスライドでは、モビリティと物流における不確実性 を通じて 不確実性は、需要・リソース・環境の変化として整理できる。 ことを理解する必要がある。試験答案では、まず概念の定義を一文で書き、次に講義の例へ対応付け、最後に意思決定モデルにどのような影響を与えるかを述べる。単語の列挙だけでは不十分であり、データがどのように情報モデルへ変換され、その情報モデルがどのように決定変数・目的関数・制約へ接続されるかまで説明する。
\item \textbf{Slide 005「モビリティと物流における不確実性」を、配送またはTSPの具体例で説明せよ。}\\注文がいつ来るか、ドライバーが稼働するか、道路が混むかは意思決定時点では確定しない。これが不確実性である。 答案では、例を出した後に、それがどのモデル要素に対応するかを明示する。例えば、顧客集合や旅行時間行列は情報モデル、訪問順序や車両割当は意思決定モデル、実際の配送指示は実装である。
\end{enumerate}
\end{minipage}
\end{minipage}
\end{adjustbox}

\clearpage
\SlideHeader{Lecture 3: Uncertainty}{006}{旅行時間の不確実性}
\SlideImage{6}{../Materials/2026/3DM_03_Uncertainty_SS26.pdf}
\begin{adjustbox}{max totalsize={\textwidth}{0.675\textheight},center}
\begin{minipage}{\textwidth}\scriptsize
\begin{minipage}[t]{0.485\textwidth}
\NoteSection{日本語訳}
\begin{itemize}
\item 天候
\begin{itemize}
\item 連続的
\item 予測可能
\end{itemize}
\item 交通
\begin{itemize}
\item システム的で観測可能なパターンに従う。
\item 事故や混雑など突発的。
\item 部分的に予測可能。
\end{itemize}
\item 交通管理
\begin{itemize}
\item 離散イベント。
\item 枠組み条件と過去イベントに依存する。
\item 部分的に予測可能。
\end{itemize}
\end{itemize}
\NoteSection{試験対策ポイント}
\begin{itemize}
\item 不確実性は、需要・リソース・環境の変化として整理できる。
\item 頻度、範囲、破壊力、予測可能性で性質が異なる。
\end{itemize}
\end{minipage}\hfill
\begin{minipage}[t]{0.485\textwidth}
\NoteSection{解説}
このページでは「旅行時間の不確実性」を、講義全体の流れの中で位置づける。スライドの箇条書きは短いが、試験ではその背後にある概念、現実例、モデル上の意味を文章で説明できる必要がある。\par
不確実性とは、意思決定時点では将来の状態や入力値が完全には分からないことである。配送であれば、注文がどこから来るか、交通がどれだけ混むか、サービス時間が何分か、ドライバーが利用可能かどうかが事前には確定していない。\par
不確実性は一種類ではない。需要不確実性は、顧客からの注文の発生時刻、場所、量に関する不確実性である。リソース不確実性は、車両故障、ドライバーのキャンセル、バッテリー残量、積載能力に関する不確実性である。環境不確実性は、交通、天候、道路工事、事故、駐車可能性に関する不確実性である。\par
不確実性は、頻度、範囲、破壊力、予測可能性によって性質が異なる。例えば日々の交通変動は頻繁に起きるが、過去データからある程度予測できる。一方で、大規模事故や道路閉鎖は低頻度だが、発生すれば計画を大きく壊す。\par
不確実性を無視して平均値だけで計画すると、平均的には良く見えるが実行時に頻繁に遅れる計画になる可能性がある。特に複数顧客を連続して訪問するルーティングでは、一つの遅延が後続訪問に波及するため、単独の平均時間以上にリスクが大きい。\par
試験答案では、不確実性の例を挙げるだけでなく、それが情報モデルと意思決定モデルにどう影響するかを書くとよい。情報モデルでは確率分布や区間として表現し、意思決定モデルでは期待値、リスク、ロバスト性、再最適化などを考える。\par
\NoteSection{確認問題と解答}
\begin{enumerate}
\item \textbf{Slide 006「旅行時間の不確実性」の中心概念を、専門用語を使わずに説明せよ。}\\このスライドでは、旅行時間の不確実性 を通じて 不確実性は、需要・リソース・環境の変化として整理できる。 ことを理解する必要がある。試験答案では、まず概念の定義を一文で書き、次に講義の例へ対応付け、最後に意思決定モデルにどのような影響を与えるかを述べる。単語の列挙だけでは不十分であり、データがどのように情報モデルへ変換され、その情報モデルがどのように決定変数・目的関数・制約へ接続されるかまで説明する。
\item \textbf{Slide 006「旅行時間の不確実性」を、配送またはTSPの具体例で説明せよ。}\\注文がいつ来るか、ドライバーが稼働するか、道路が混むかは意思決定時点では確定しない。これが不確実性である。 答案では、例を出した後に、それがどのモデル要素に対応するかを明示する。例えば、顧客集合や旅行時間行列は情報モデル、訪問順序や車両割当は意思決定モデル、実際の配送指示は実装である。
\end{enumerate}
\end{minipage}
\end{minipage}
\end{adjustbox}

\clearpage
\SlideHeader{Lecture 3: Uncertainty}{007}{起こり得る結果}
\SlideImage{7}{../Materials/2026/3DM_03_Uncertainty_SS26.pdf}
\begin{adjustbox}{max totalsize={\textwidth}{0.675\textheight},center}
\begin{minipage}{\textwidth}\scriptsize
\begin{minipage}[t]{0.485\textwidth}
\NoteSection{日本語訳}
\begin{itemize}
\item 計画活動の完全なキャンセル
\begin{itemize}
\item 輸送: 荷物が配送されない。
\item モビリティ: 職場に到達できない。
\end{itemize}
\item 計画活動の中断
\begin{itemize}
\item 輸送: 荷物配送が遅れる。顧客が不在になった場合、完全な失敗になり得る。
\item モビリティ: 職場に遅れて到着する。
\end{itemize}
\item 一般的に望ましくない効果
\begin{itemize}
\item 追加費用が発生する。
\item 顧客満足度が低下する。
\item 後続活動が影響を受ける。
\end{itemize}
\item 意思決定において不確実性を考慮することは重要である。
\end{itemize}
\NoteSection{試験対策ポイント}
\begin{itemize}
\item 不確実性は、需要・リソース・環境の変化として整理できる。
\item 頻度、範囲、破壊力、予測可能性で性質が異なる。
\end{itemize}
\end{minipage}\hfill
\begin{minipage}[t]{0.485\textwidth}
\NoteSection{解説}
このページでは「起こり得る結果」を、講義全体の流れの中で位置づける。スライドの箇条書きは短いが、試験ではその背後にある概念、現実例、モデル上の意味を文章で説明できる必要がある。\par
不確実性とは、意思決定時点では将来の状態や入力値が完全には分からないことである。配送であれば、注文がどこから来るか、交通がどれだけ混むか、サービス時間が何分か、ドライバーが利用可能かどうかが事前には確定していない。\par
不確実性は一種類ではない。需要不確実性は、顧客からの注文の発生時刻、場所、量に関する不確実性である。リソース不確実性は、車両故障、ドライバーのキャンセル、バッテリー残量、積載能力に関する不確実性である。環境不確実性は、交通、天候、道路工事、事故、駐車可能性に関する不確実性である。\par
不確実性は、頻度、範囲、破壊力、予測可能性によって性質が異なる。例えば日々の交通変動は頻繁に起きるが、過去データからある程度予測できる。一方で、大規模事故や道路閉鎖は低頻度だが、発生すれば計画を大きく壊す。\par
不確実性を無視して平均値だけで計画すると、平均的には良く見えるが実行時に頻繁に遅れる計画になる可能性がある。特に複数顧客を連続して訪問するルーティングでは、一つの遅延が後続訪問に波及するため、単独の平均時間以上にリスクが大きい。\par
試験答案では、不確実性の例を挙げるだけでなく、それが情報モデルと意思決定モデルにどう影響するかを書くとよい。情報モデルでは確率分布や区間として表現し、意思決定モデルでは期待値、リスク、ロバスト性、再最適化などを考える。\par
\NoteSection{確認問題と解答}
\begin{enumerate}
\item \textbf{Slide 007「起こり得る結果」の中心概念を、専門用語を使わずに説明せよ。}\\このスライドでは、起こり得る結果 を通じて 不確実性は、需要・リソース・環境の変化として整理できる。 ことを理解する必要がある。試験答案では、まず概念の定義を一文で書き、次に講義の例へ対応付け、最後に意思決定モデルにどのような影響を与えるかを述べる。単語の列挙だけでは不十分であり、データがどのように情報モデルへ変換され、その情報モデルがどのように決定変数・目的関数・制約へ接続されるかまで説明する。
\item \textbf{Slide 007「起こり得る結果」を、配送またはTSPの具体例で説明せよ。}\\注文がいつ来るか、ドライバーが稼働するか、道路が混むかは意思決定時点では確定しない。これが不確実性である。 答案では、例を出した後に、それがどのモデル要素に対応するかを明示する。例えば、顧客集合や旅行時間行列は情報モデル、訪問順序や車両割当は意思決定モデル、実際の配送指示は実装である。
\end{enumerate}
\end{minipage}
\end{minipage}
\end{adjustbox}

\clearpage
\SlideHeader{Lecture 3: Uncertainty}{008}{アジェンダ}
\SlideImage{8}{../Materials/2026/3DM_03_Uncertainty_SS26.pdf}
\begin{adjustbox}{max totalsize={\textwidth}{0.675\textheight},center}
\begin{minipage}{\textwidth}\scriptsize
\begin{minipage}[t]{0.485\textwidth}
\NoteSection{日本語訳}
\begin{itemize}
\item 不確実性。
\item 不確実性のモデリング。
\item 不確実性下の意思決定。
\item ケーススタディ: 都市物流におけるロバスト車両ルーティング。
\end{itemize}
\end{minipage}\hfill
\begin{minipage}[t]{0.485\textwidth}
\NoteSection{注記}
{\small このページは表紙・目次・連絡先などの事務情報であるため、解説と確認問題は付けない。}
\end{minipage}
\end{minipage}
\end{adjustbox}

\clearpage
\SlideHeader{Lecture 3: Uncertainty}{009}{不確実性のモデリング}
\SlideImage{9}{../Materials/2026/3DM_03_Uncertainty_SS26.pdf}
\begin{adjustbox}{max totalsize={\textwidth}{0.675\textheight},center}
\begin{minipage}{\textwidth}\scriptsize
\begin{minipage}[t]{0.485\textwidth}
\NoteSection{日本語訳}
\begin{itemize}
\item モデルの選択は情報レベルに依存する。
\item 1. Deterministic: すべての情報が既知である。
\item 2. Stochastic: 異なる情報実現の確率が既知である。
\item 3. Quasi-stochastic: 情報実現の確率が未知である。
\end{itemize}
\NoteSection{試験対策ポイント}
\begin{itemize}
\item deterministic, stochastic, quasi-stochastic の違いは、将来値と確率情報をどう扱うかで決まる。
\item 確率分布がある場合とない場合で、使える評価基準が変わる。
\end{itemize}
\end{minipage}\hfill
\begin{minipage}[t]{0.485\textwidth}
\NoteSection{解説}
このページでは「不確実性のモデリング」を、講義全体の流れの中で位置づける。スライドの箇条書きは短いが、試験ではその背後にある概念、現実例、モデル上の意味を文章で説明できる必要がある。\par
Deterministic model は、入力値が既知で固定されていると仮定する。例えば、AからBへの旅行時間を常に15分と置く。これは単純で計算しやすいが、実際の交通変動やサービス時間のばらつきを無視するため、現実で計画が崩れることがある。\par
Stochastic model は、不確実な変数を確率分布として表す。旅行時間が平均20分・標準偏差5分の分布に従う、注文発生がポアソン過程に従う、サービス時間が対数正規分布に従う、といった表現である。確率が分かるため、期待費用、分散、分位点、確率制約、サンプリングを使える。\par
Quasi-stochastic model は、確率分布は分からないが、起こり得るシナリオや値の範囲は分かる場合に使う。例えば、旅行時間が10分から25分の間であることは分かるが、その確率は分からない場合である。このとき、区間、代表シナリオ、最悪ケース、Hurwicz基準、minmax regret が使われる。\par
確率分布が十分に推定できるなら stochastic model が自然である。しかし、過去データが少ない、新しいサービスで履歴がない、環境が変化して過去データが当てにならない場合は quasi-stochastic model が有用である。どちらを使うかは、データ量、分布仮定の妥当性、意思決定の目的によって決まる。\par
試験では、stochastic と quasi-stochastic の違いを『不確実であるかどうか』ではなく『確率情報を持っているかどうか』で説明する。どちらも不確実性を扱うが、stochastic は確率を使い、quasi-stochastic は確率なしで区間やシナリオを使う。\par
\NoteSection{確認問題と解答}
\begin{enumerate}
\item \textbf{Slide 009「不確実性のモデリング」の中心概念を、専門用語を使わずに説明せよ。}\\このスライドでは、不確実性のモデリング を通じて deterministic, stochastic, quasi-stochastic の違いは、将来値と確率情報をどう扱うかで決まる。 ことを理解する必要がある。試験答案では、まず概念の定義を一文で書き、次に講義の例へ対応付け、最後に意思決定モデルにどのような影響を与えるかを述べる。単語の列挙だけでは不十分であり、データがどのように情報モデルへ変換され、その情報モデルがどのように決定変数・目的関数・制約へ接続されるかまで説明する。
\item \textbf{Slide 009「不確実性のモデリング」を、配送またはTSPの具体例で説明せよ。}\\旅行時間分布が分かるならstochastic、10-25分という範囲だけ分かるならquasi-stochasticとして扱う。 答案では、例を出した後に、それがどのモデル要素に対応するかを明示する。例えば、顧客集合や旅行時間行列は情報モデル、訪問順序や車両割当は意思決定モデル、実際の配送指示は実装である。
\end{enumerate}
\end{minipage}
\end{minipage}
\end{adjustbox}

\clearpage
\SlideHeader{Lecture 3: Uncertainty}{010}{基礎}
\SlideImage{10}{../Materials/2026/3DM_03_Uncertainty_SS26.pdf}
\begin{adjustbox}{max totalsize={\textwidth}{0.675\textheight},center}
\begin{minipage}{\textwidth}\scriptsize
\begin{minipage}[t]{0.485\textwidth}
\NoteSection{日本語訳}
\begin{itemize}
\item 確率分布。
\item 離散または連続。
\item 確率変数 X は実現値 x\_1, ..., x\_i を持つ。例: x\_1=20 km/h。
\item 対応する実現値 x\_i の確率 p\_1, ..., p\_i。
\item 期待値: mu = E[X] = sum\_i p\_i x\_i、または積分。
\item 分散。
\item 標準偏差。
\item 変動係数: CV = sigma / mu。
\item 中央値と分位点。
\item これらの指標は意思決定プロセスに組み込める。
\item Sampling: 分布から新しい観測を生成すること。
\end{itemize}
\NoteSection{試験対策ポイント}
\begin{itemize}
\item 期待値は平均性能、分散や分位点はリスクを表す。
\item 分布フィットでは可視化、パラメータ推定、適合度確認が必要である。
\end{itemize}
\end{minipage}\hfill
\begin{minipage}[t]{0.485\textwidth}
\NoteSection{解説}
このページでは「基礎」を、講義全体の流れの中で位置づける。スライドの箇条書きは短いが、試験ではその背後にある概念、現実例、モデル上の意味を文章で説明できる必要がある。\par
確率分布は、不確実な変数がどの値をどれくらいの確率で取るかを表す。配送では、旅行時間、サービス時間、注文間隔、需要量、キャンセル数などを分布で表せる。分布を使うことで、単一の平均値ではなく、ばらつきや極端な値も考慮できる。\par
期待値は長期的な平均を表すが、期待値だけでは不十分なことが多い。平均20分の旅行時間でも、ほとんどが18-22分なのか、10分と50分が混在しているのかでは、計画上のリスクが全く違う。分散や標準偏差は、このばらつきの大きさを表す。\par
分位点はサービスレベルを考える際に重要である。例えば90\%分位点が35分なら、90\%のケースで35分以内に到着することを意味する。顧客に到着時間を約束する場合、平均よりも上位分位点を見る方が安全な計画につながる。\par
実データから分布を決めるには、まず観測値を集め、外れ値や欠損を確認し、ヒストグラムや箱ひげ図で形を見る。その後、正規分布、対数正規分布、ガンマ分布など候補を選び、平均や分散などのパラメータを推定し、適合度を確認する。\par
サービス時間のように負にならない値では、正規分布よりも右裾を持つ非負分布が自然な場合がある。家電修理や在宅介護では、多くの作業は短時間で終わるが、例外的に長い作業が起こる。試験では、分布名だけでなく、なぜその形が現実に合うのかを説明する。\par
\NoteSection{確認問題と解答}
\begin{enumerate}
\item \textbf{Slide 010「基礎」の中心概念を、専門用語を使わずに説明せよ。}\\このスライドでは、基礎 を通じて 期待値は平均性能、分散や分位点はリスクを表す。 ことを理解する必要がある。試験答案では、まず概念の定義を一文で書き、次に講義の例へ対応付け、最後に意思決定モデルにどのような影響を与えるかを述べる。単語の列挙だけでは不十分であり、データがどのように情報モデルへ変換され、その情報モデルがどのように決定変数・目的関数・制約へ接続されるかまで説明する。
\item \textbf{Slide 010「基礎」を、配送またはTSPの具体例で説明せよ。}\\旅行時間の平均だけでなく、標準偏差や90\%分位点を見ることで遅延リスクを評価する。 答案では、例を出した後に、それがどのモデル要素に対応するかを明示する。例えば、顧客集合や旅行時間行列は情報モデル、訪問順序や車両割当は意思決定モデル、実際の配送指示は実装である。
\end{enumerate}
\end{minipage}
\end{minipage}
\end{adjustbox}

\clearpage
\SlideHeader{Lecture 3: Uncertainty}{011}{確率分布の決定}
\SlideImage{11}{../Materials/2026/3DM_03_Uncertainty_SS26.pdf}
\begin{adjustbox}{max totalsize={\textwidth}{0.675\textheight},center}
\begin{minipage}{\textwidth}\scriptsize
\begin{minipage}[t]{0.485\textwidth}
\NoteSection{日本語訳}
\begin{itemize}
\item 理論分布と対応パラメータを単純に仮定する。
\item 例
\begin{itemize}
\item 一様分布
\item 正規分布
\item 指数分布
\end{itemize}
\item 既存データへの確率分布の選択と調整、つまり fitting。
\item 手順
\begin{itemize}
\item 適切な理論分布を選ぶ。
\item 対応するパラメータを推定する。
\item 適合度を検査する。可視化、Kolmogorov-Smirnov検定など。
\item 推定パラメータの妥当性をブートストラップで検査する。
\end{itemize}
\end{itemize}
\NoteSection{試験対策ポイント}
\begin{itemize}
\item 期待値は平均性能、分散や分位点はリスクを表す。
\item 分布フィットでは可視化、パラメータ推定、適合度確認が必要である。
\end{itemize}
\end{minipage}\hfill
\begin{minipage}[t]{0.485\textwidth}
\NoteSection{解説}
このページでは「確率分布の決定」を、講義全体の流れの中で位置づける。スライドの箇条書きは短いが、試験ではその背後にある概念、現実例、モデル上の意味を文章で説明できる必要がある。\par
確率分布は、不確実な変数がどの値をどれくらいの確率で取るかを表す。配送では、旅行時間、サービス時間、注文間隔、需要量、キャンセル数などを分布で表せる。分布を使うことで、単一の平均値ではなく、ばらつきや極端な値も考慮できる。\par
期待値は長期的な平均を表すが、期待値だけでは不十分なことが多い。平均20分の旅行時間でも、ほとんどが18-22分なのか、10分と50分が混在しているのかでは、計画上のリスクが全く違う。分散や標準偏差は、このばらつきの大きさを表す。\par
分位点はサービスレベルを考える際に重要である。例えば90\%分位点が35分なら、90\%のケースで35分以内に到着することを意味する。顧客に到着時間を約束する場合、平均よりも上位分位点を見る方が安全な計画につながる。\par
実データから分布を決めるには、まず観測値を集め、外れ値や欠損を確認し、ヒストグラムや箱ひげ図で形を見る。その後、正規分布、対数正規分布、ガンマ分布など候補を選び、平均や分散などのパラメータを推定し、適合度を確認する。\par
サービス時間のように負にならない値では、正規分布よりも右裾を持つ非負分布が自然な場合がある。家電修理や在宅介護では、多くの作業は短時間で終わるが、例外的に長い作業が起こる。試験では、分布名だけでなく、なぜその形が現実に合うのかを説明する。\par
\NoteSection{確認問題と解答}
\begin{enumerate}
\item \textbf{Slide 011「確率分布の決定」の中心概念を、専門用語を使わずに説明せよ。}\\このスライドでは、確率分布の決定 を通じて 期待値は平均性能、分散や分位点はリスクを表す。 ことを理解する必要がある。試験答案では、まず概念の定義を一文で書き、次に講義の例へ対応付け、最後に意思決定モデルにどのような影響を与えるかを述べる。単語の列挙だけでは不十分であり、データがどのように情報モデルへ変換され、その情報モデルがどのように決定変数・目的関数・制約へ接続されるかまで説明する。
\item \textbf{Slide 011「確率分布の決定」を、配送またはTSPの具体例で説明せよ。}\\旅行時間の平均だけでなく、標準偏差や90\%分位点を見ることで遅延リスクを評価する。 答案では、例を出した後に、それがどのモデル要素に対応するかを明示する。例えば、顧客集合や旅行時間行列は情報モデル、訪問順序や車両割当は意思決定モデル、実際の配送指示は実装である。
\end{enumerate}
\end{minipage}
\end{minipage}
\end{adjustbox}

\clearpage
\SlideHeader{Lecture 3: Uncertainty}{012}{理論的な確率密度関数}
\SlideImage{12}{../Materials/2026/3DM_03_Uncertainty_SS26.pdf}
\begin{adjustbox}{max totalsize={\textwidth}{0.675\textheight},center}
\begin{minipage}{\textwidth}\scriptsize
\begin{minipage}[t]{0.485\textwidth}
\NoteSection{日本語訳}
\begin{itemize}
\item 複数の理論的確率密度関数の形。
\end{itemize}
\NoteSection{試験対策ポイント}
\begin{itemize}
\item 期待値は平均性能、分散や分位点はリスクを表す。
\item 分布フィットでは可視化、パラメータ推定、適合度確認が必要である。
\end{itemize}
\end{minipage}\hfill
\begin{minipage}[t]{0.485\textwidth}
\NoteSection{解説}
このページでは「理論的な確率密度関数」を、講義全体の流れの中で位置づける。スライドの箇条書きは短いが、試験ではその背後にある概念、現実例、モデル上の意味を文章で説明できる必要がある。\par
確率分布は、不確実な変数がどの値をどれくらいの確率で取るかを表す。配送では、旅行時間、サービス時間、注文間隔、需要量、キャンセル数などを分布で表せる。分布を使うことで、単一の平均値ではなく、ばらつきや極端な値も考慮できる。\par
期待値は長期的な平均を表すが、期待値だけでは不十分なことが多い。平均20分の旅行時間でも、ほとんどが18-22分なのか、10分と50分が混在しているのかでは、計画上のリスクが全く違う。分散や標準偏差は、このばらつきの大きさを表す。\par
分位点はサービスレベルを考える際に重要である。例えば90\%分位点が35分なら、90\%のケースで35分以内に到着することを意味する。顧客に到着時間を約束する場合、平均よりも上位分位点を見る方が安全な計画につながる。\par
実データから分布を決めるには、まず観測値を集め、外れ値や欠損を確認し、ヒストグラムや箱ひげ図で形を見る。その後、正規分布、対数正規分布、ガンマ分布など候補を選び、平均や分散などのパラメータを推定し、適合度を確認する。\par
サービス時間のように負にならない値では、正規分布よりも右裾を持つ非負分布が自然な場合がある。家電修理や在宅介護では、多くの作業は短時間で終わるが、例外的に長い作業が起こる。試験では、分布名だけでなく、なぜその形が現実に合うのかを説明する。\par
\NoteSection{確認問題と解答}
\begin{enumerate}
\item \textbf{Slide 012「理論的な確率密度関数」の中心概念を、専門用語を使わずに説明せよ。}\\このスライドでは、理論的な確率密度関数 を通じて 期待値は平均性能、分散や分位点はリスクを表す。 ことを理解する必要がある。試験答案では、まず概念の定義を一文で書き、次に講義の例へ対応付け、最後に意思決定モデルにどのような影響を与えるかを述べる。単語の列挙だけでは不十分であり、データがどのように情報モデルへ変換され、その情報モデルがどのように決定変数・目的関数・制約へ接続されるかまで説明する。
\item \textbf{Slide 012「理論的な確率密度関数」を、配送またはTSPの具体例で説明せよ。}\\旅行時間の平均だけでなく、標準偏差や90\%分位点を見ることで遅延リスクを評価する。 答案では、例を出した後に、それがどのモデル要素に対応するかを明示する。例えば、顧客集合や旅行時間行列は情報モデル、訪問順序や車両割当は意思決定モデル、実際の配送指示は実装である。
\end{enumerate}
\end{minipage}
\end{minipage}
\end{adjustbox}

\clearpage
\SlideHeader{Lecture 3: Uncertainty}{013}{例: 都市旅行時間分布 I}
\SlideImage{13}{../Materials/2026/3DM_03_Uncertainty_SS26.pdf}
\begin{adjustbox}{max totalsize={\textwidth}{0.675\textheight},center}
\begin{minipage}{\textwidth}\scriptsize
\begin{minipage}[t]{0.485\textwidth}
\NoteSection{日本語訳}
\begin{itemize}
\item 一つの道路区間における旅行時間観測。
\item 横軸は旅行時間、縦軸は頻度。
\end{itemize}
\NoteSection{試験対策ポイント}
\begin{itemize}
\item 期待値は平均性能、分散や分位点はリスクを表す。
\item 分布フィットでは可視化、パラメータ推定、適合度確認が必要である。
\end{itemize}
\end{minipage}\hfill
\begin{minipage}[t]{0.485\textwidth}
\NoteSection{解説}
このページでは「例: 都市旅行時間分布 I」を、講義全体の流れの中で位置づける。スライドの箇条書きは短いが、試験ではその背後にある概念、現実例、モデル上の意味を文章で説明できる必要がある。\par
確率分布は、不確実な変数がどの値をどれくらいの確率で取るかを表す。配送では、旅行時間、サービス時間、注文間隔、需要量、キャンセル数などを分布で表せる。分布を使うことで、単一の平均値ではなく、ばらつきや極端な値も考慮できる。\par
期待値は長期的な平均を表すが、期待値だけでは不十分なことが多い。平均20分の旅行時間でも、ほとんどが18-22分なのか、10分と50分が混在しているのかでは、計画上のリスクが全く違う。分散や標準偏差は、このばらつきの大きさを表す。\par
分位点はサービスレベルを考える際に重要である。例えば90\%分位点が35分なら、90\%のケースで35分以内に到着することを意味する。顧客に到着時間を約束する場合、平均よりも上位分位点を見る方が安全な計画につながる。\par
実データから分布を決めるには、まず観測値を集め、外れ値や欠損を確認し、ヒストグラムや箱ひげ図で形を見る。その後、正規分布、対数正規分布、ガンマ分布など候補を選び、平均や分散などのパラメータを推定し、適合度を確認する。\par
サービス時間のように負にならない値では、正規分布よりも右裾を持つ非負分布が自然な場合がある。家電修理や在宅介護では、多くの作業は短時間で終わるが、例外的に長い作業が起こる。試験では、分布名だけでなく、なぜその形が現実に合うのかを説明する。\par
\NoteSection{確認問題と解答}
\begin{enumerate}
\item \textbf{Slide 013「例: 都市旅行時間分布 I」の中心概念を、専門用語を使わずに説明せよ。}\\このスライドでは、例: 都市旅行時間分布 I を通じて 期待値は平均性能、分散や分位点はリスクを表す。 ことを理解する必要がある。試験答案では、まず概念の定義を一文で書き、次に講義の例へ対応付け、最後に意思決定モデルにどのような影響を与えるかを述べる。単語の列挙だけでは不十分であり、データがどのように情報モデルへ変換され、その情報モデルがどのように決定変数・目的関数・制約へ接続されるかまで説明する。
\item \textbf{Slide 013「例: 都市旅行時間分布 I」を、配送またはTSPの具体例で説明せよ。}\\旅行時間の平均だけでなく、標準偏差や90\%分位点を見ることで遅延リスクを評価する。 答案では、例を出した後に、それがどのモデル要素に対応するかを明示する。例えば、顧客集合や旅行時間行列は情報モデル、訪問順序や車両割当は意思決定モデル、実際の配送指示は実装である。
\end{enumerate}
\end{minipage}
\end{minipage}
\end{adjustbox}

\clearpage
\SlideHeader{Lecture 3: Uncertainty}{014}{ソフトウェアによるフィッティング}
\SlideImage{14}{../Materials/2026/3DM_03_Uncertainty_SS26.pdf}
\begin{adjustbox}{max totalsize={\textwidth}{0.675\textheight},center}
\begin{minipage}{\textwidth}\scriptsize
\begin{minipage}[t]{0.485\textwidth}
\NoteSection{日本語訳}
\begin{itemize}
\item 支援ツール: Python とライブラリ。
\item 自動ツール: EasyFit。
\end{itemize}
\NoteSection{試験対策ポイント}
\begin{itemize}
\item 期待値は平均性能、分散や分位点はリスクを表す。
\item 分布フィットでは可視化、パラメータ推定、適合度確認が必要である。
\end{itemize}
\end{minipage}\hfill
\begin{minipage}[t]{0.485\textwidth}
\NoteSection{解説}
このページでは「ソフトウェアによるフィッティング」を、講義全体の流れの中で位置づける。スライドの箇条書きは短いが、試験ではその背後にある概念、現実例、モデル上の意味を文章で説明できる必要がある。\par
確率分布は、不確実な変数がどの値をどれくらいの確率で取るかを表す。配送では、旅行時間、サービス時間、注文間隔、需要量、キャンセル数などを分布で表せる。分布を使うことで、単一の平均値ではなく、ばらつきや極端な値も考慮できる。\par
期待値は長期的な平均を表すが、期待値だけでは不十分なことが多い。平均20分の旅行時間でも、ほとんどが18-22分なのか、10分と50分が混在しているのかでは、計画上のリスクが全く違う。分散や標準偏差は、このばらつきの大きさを表す。\par
分位点はサービスレベルを考える際に重要である。例えば90\%分位点が35分なら、90\%のケースで35分以内に到着することを意味する。顧客に到着時間を約束する場合、平均よりも上位分位点を見る方が安全な計画につながる。\par
実データから分布を決めるには、まず観測値を集め、外れ値や欠損を確認し、ヒストグラムや箱ひげ図で形を見る。その後、正規分布、対数正規分布、ガンマ分布など候補を選び、平均や分散などのパラメータを推定し、適合度を確認する。\par
サービス時間のように負にならない値では、正規分布よりも右裾を持つ非負分布が自然な場合がある。家電修理や在宅介護では、多くの作業は短時間で終わるが、例外的に長い作業が起こる。試験では、分布名だけでなく、なぜその形が現実に合うのかを説明する。\par
\NoteSection{確認問題と解答}
\begin{enumerate}
\item \textbf{Slide 014「ソフトウェアによるフィッティング」の中心概念を、専門用語を使わずに説明せよ。}\\このスライドでは、ソフトウェアによるフィッティング を通じて 期待値は平均性能、分散や分位点はリスクを表す。 ことを理解する必要がある。試験答案では、まず概念の定義を一文で書き、次に講義の例へ対応付け、最後に意思決定モデルにどのような影響を与えるかを述べる。単語の列挙だけでは不十分であり、データがどのように情報モデルへ変換され、その情報モデルがどのように決定変数・目的関数・制約へ接続されるかまで説明する。
\item \textbf{Slide 014「ソフトウェアによるフィッティング」を、配送またはTSPの具体例で説明せよ。}\\旅行時間の平均だけでなく、標準偏差や90\%分位点を見ることで遅延リスクを評価する。 答案では、例を出した後に、それがどのモデル要素に対応するかを明示する。例えば、顧客集合や旅行時間行列は情報モデル、訪問順序や車両割当は意思決定モデル、実際の配送指示は実装である。
\end{enumerate}
\end{minipage}
\end{minipage}
\end{adjustbox}

\clearpage
\SlideHeader{Lecture 3: Uncertainty}{015}{例: 都市旅行時間分布 II}
\SlideImage{15}{../Materials/2026/3DM_03_Uncertainty_SS26.pdf}
\begin{adjustbox}{max totalsize={\textwidth}{0.675\textheight},center}
\begin{minipage}{\textwidth}\scriptsize
\begin{minipage}[t]{0.485\textwidth}
\NoteSection{日本語訳}
\begin{itemize}
\item 旅行時間の頻度または密度。
\item 観測分布に理論分布を当てはめる。
\end{itemize}
\NoteSection{試験対策ポイント}
\begin{itemize}
\item 期待値は平均性能、分散や分位点はリスクを表す。
\item 分布フィットでは可視化、パラメータ推定、適合度確認が必要である。
\end{itemize}
\end{minipage}\hfill
\begin{minipage}[t]{0.485\textwidth}
\NoteSection{解説}
このページでは「例: 都市旅行時間分布 II」を、講義全体の流れの中で位置づける。スライドの箇条書きは短いが、試験ではその背後にある概念、現実例、モデル上の意味を文章で説明できる必要がある。\par
確率分布は、不確実な変数がどの値をどれくらいの確率で取るかを表す。配送では、旅行時間、サービス時間、注文間隔、需要量、キャンセル数などを分布で表せる。分布を使うことで、単一の平均値ではなく、ばらつきや極端な値も考慮できる。\par
期待値は長期的な平均を表すが、期待値だけでは不十分なことが多い。平均20分の旅行時間でも、ほとんどが18-22分なのか、10分と50分が混在しているのかでは、計画上のリスクが全く違う。分散や標準偏差は、このばらつきの大きさを表す。\par
分位点はサービスレベルを考える際に重要である。例えば90\%分位点が35分なら、90\%のケースで35分以内に到着することを意味する。顧客に到着時間を約束する場合、平均よりも上位分位点を見る方が安全な計画につながる。\par
実データから分布を決めるには、まず観測値を集め、外れ値や欠損を確認し、ヒストグラムや箱ひげ図で形を見る。その後、正規分布、対数正規分布、ガンマ分布など候補を選び、平均や分散などのパラメータを推定し、適合度を確認する。\par
サービス時間のように負にならない値では、正規分布よりも右裾を持つ非負分布が自然な場合がある。家電修理や在宅介護では、多くの作業は短時間で終わるが、例外的に長い作業が起こる。試験では、分布名だけでなく、なぜその形が現実に合うのかを説明する。\par
\NoteSection{確認問題と解答}
\begin{enumerate}
\item \textbf{Slide 015「例: 都市旅行時間分布 II」の中心概念を、専門用語を使わずに説明せよ。}\\このスライドでは、例: 都市旅行時間分布 II を通じて 期待値は平均性能、分散や分位点はリスクを表す。 ことを理解する必要がある。試験答案では、まず概念の定義を一文で書き、次に講義の例へ対応付け、最後に意思決定モデルにどのような影響を与えるかを述べる。単語の列挙だけでは不十分であり、データがどのように情報モデルへ変換され、その情報モデルがどのように決定変数・目的関数・制約へ接続されるかまで説明する。
\item \textbf{Slide 015「例: 都市旅行時間分布 II」を、配送またはTSPの具体例で説明せよ。}\\旅行時間の平均だけでなく、標準偏差や90\%分位点を見ることで遅延リスクを評価する。 答案では、例を出した後に、それがどのモデル要素に対応するかを明示する。例えば、顧客集合や旅行時間行列は情報モデル、訪問順序や車両割当は意思決定モデル、実際の配送指示は実装である。
\end{enumerate}
\end{minipage}
\end{minipage}
\end{adjustbox}

\clearpage
\SlideHeader{Lecture 3: Uncertainty}{016}{確率的モデリングの問題点}
\SlideImage{16}{../Materials/2026/3DM_03_Uncertainty_SS26.pdf}
\begin{adjustbox}{max totalsize={\textwidth}{0.675\textheight},center}
\begin{minipage}{\textwidth}\scriptsize
\begin{minipage}[t]{0.485\textwidth}
\NoteSection{日本語訳}
\begin{itemize}
\item データが利用できない
\begin{itemize}
\item 事前に観測できない。
\item データ記録がない。
\item 道路工事などシステムが変化する。
\item 自動運転車など新しい概念で履歴がない。
\end{itemize}
\item 利用可能データが限られる
\begin{itemize}
\item 観測が少なく、基礎分布の推定が粗くなる。
\end{itemize}
\item 道路セグメントごとの独立分布を畳み込んで経路旅行時間を計算する仮定では、相関が無視される。
\item 実際には、隣接道路セグメントの相関が重要である。
\end{itemize}
\NoteSection{試験対策ポイント}
\begin{itemize}
\item deterministic, stochastic, quasi-stochastic の違いは、将来値と確率情報をどう扱うかで決まる。
\item 確率分布がある場合とない場合で、使える評価基準が変わる。
\end{itemize}
\end{minipage}\hfill
\begin{minipage}[t]{0.485\textwidth}
\NoteSection{解説}
このページでは「確率的モデリングの問題点」を、講義全体の流れの中で位置づける。スライドの箇条書きは短いが、試験ではその背後にある概念、現実例、モデル上の意味を文章で説明できる必要がある。\par
Deterministic model は、入力値が既知で固定されていると仮定する。例えば、AからBへの旅行時間を常に15分と置く。これは単純で計算しやすいが、実際の交通変動やサービス時間のばらつきを無視するため、現実で計画が崩れることがある。\par
Stochastic model は、不確実な変数を確率分布として表す。旅行時間が平均20分・標準偏差5分の分布に従う、注文発生がポアソン過程に従う、サービス時間が対数正規分布に従う、といった表現である。確率が分かるため、期待費用、分散、分位点、確率制約、サンプリングを使える。\par
Quasi-stochastic model は、確率分布は分からないが、起こり得るシナリオや値の範囲は分かる場合に使う。例えば、旅行時間が10分から25分の間であることは分かるが、その確率は分からない場合である。このとき、区間、代表シナリオ、最悪ケース、Hurwicz基準、minmax regret が使われる。\par
確率分布が十分に推定できるなら stochastic model が自然である。しかし、過去データが少ない、新しいサービスで履歴がない、環境が変化して過去データが当てにならない場合は quasi-stochastic model が有用である。どちらを使うかは、データ量、分布仮定の妥当性、意思決定の目的によって決まる。\par
試験では、stochastic と quasi-stochastic の違いを『不確実であるかどうか』ではなく『確率情報を持っているかどうか』で説明する。どちらも不確実性を扱うが、stochastic は確率を使い、quasi-stochastic は確率なしで区間やシナリオを使う。\par
\NoteSection{確認問題と解答}
\begin{enumerate}
\item \textbf{Slide 016「確率的モデリングの問題点」の中心概念を、専門用語を使わずに説明せよ。}\\このスライドでは、確率的モデリングの問題点 を通じて deterministic, stochastic, quasi-stochastic の違いは、将来値と確率情報をどう扱うかで決まる。 ことを理解する必要がある。試験答案では、まず概念の定義を一文で書き、次に講義の例へ対応付け、最後に意思決定モデルにどのような影響を与えるかを述べる。単語の列挙だけでは不十分であり、データがどのように情報モデルへ変換され、その情報モデルがどのように決定変数・目的関数・制約へ接続されるかまで説明する。
\item \textbf{Slide 016「確率的モデリングの問題点」を、配送またはTSPの具体例で説明せよ。}\\旅行時間分布が分かるならstochastic、10-25分という範囲だけ分かるならquasi-stochasticとして扱う。 答案では、例を出した後に、それがどのモデル要素に対応するかを明示する。例えば、顧客集合や旅行時間行列は情報モデル、訪問順序や車両割当は意思決定モデル、実際の配送指示は実装である。
\end{enumerate}
\end{minipage}
\end{minipage}
\end{adjustbox}

\clearpage
\SlideHeader{Lecture 3: Uncertainty}{017}{準確率的モデリング}
\SlideImage{17}{../Materials/2026/3DM_03_Uncertainty_SS26.pdf}
\begin{adjustbox}{max totalsize={\textwidth}{0.675\textheight},center}
\begin{minipage}{\textwidth}\scriptsize
\begin{minipage}[t]{0.485\textwidth}
\NoteSection{日本語訳}
\begin{itemize}
\item シナリオに基づく
\begin{itemize}
\item n個のシナリオを作る。
\item 可能な実現値は有限個である。
\item What-if分析。
\item 例: 混雑時10 km/h、非混雑時50 km/h。
\end{itemize}
\item 区間に基づく
\begin{itemize}
\item 可能な実現値の範囲として [l,u] を使う。
\item 最小可能値 l と最大可能値 u。
\item 実現値は区間内の任意の値を取り得る。
\end{itemize}
\item シナリオと区間は、ドメイン知識またはデータから導出できる。
\item 確率分布を直接モデル化せず、不確実性の推定を表す。
\end{itemize}
\NoteSection{試験対策ポイント}
\begin{itemize}
\item deterministic, stochastic, quasi-stochastic の違いは、将来値と確率情報をどう扱うかで決まる。
\item 確率分布がある場合とない場合で、使える評価基準が変わる。
\end{itemize}
\end{minipage}\hfill
\begin{minipage}[t]{0.485\textwidth}
\NoteSection{解説}
このページでは「準確率的モデリング」を、講義全体の流れの中で位置づける。スライドの箇条書きは短いが、試験ではその背後にある概念、現実例、モデル上の意味を文章で説明できる必要がある。\par
Deterministic model は、入力値が既知で固定されていると仮定する。例えば、AからBへの旅行時間を常に15分と置く。これは単純で計算しやすいが、実際の交通変動やサービス時間のばらつきを無視するため、現実で計画が崩れることがある。\par
Stochastic model は、不確実な変数を確率分布として表す。旅行時間が平均20分・標準偏差5分の分布に従う、注文発生がポアソン過程に従う、サービス時間が対数正規分布に従う、といった表現である。確率が分かるため、期待費用、分散、分位点、確率制約、サンプリングを使える。\par
Quasi-stochastic model は、確率分布は分からないが、起こり得るシナリオや値の範囲は分かる場合に使う。例えば、旅行時間が10分から25分の間であることは分かるが、その確率は分からない場合である。このとき、区間、代表シナリオ、最悪ケース、Hurwicz基準、minmax regret が使われる。\par
確率分布が十分に推定できるなら stochastic model が自然である。しかし、過去データが少ない、新しいサービスで履歴がない、環境が変化して過去データが当てにならない場合は quasi-stochastic model が有用である。どちらを使うかは、データ量、分布仮定の妥当性、意思決定の目的によって決まる。\par
試験では、stochastic と quasi-stochastic の違いを『不確実であるかどうか』ではなく『確率情報を持っているかどうか』で説明する。どちらも不確実性を扱うが、stochastic は確率を使い、quasi-stochastic は確率なしで区間やシナリオを使う。\par
\NoteSection{確認問題と解答}
\begin{enumerate}
\item \textbf{Slide 017「準確率的モデリング」の中心概念を、専門用語を使わずに説明せよ。}\\このスライドでは、準確率的モデリング を通じて deterministic, stochastic, quasi-stochastic の違いは、将来値と確率情報をどう扱うかで決まる。 ことを理解する必要がある。試験答案では、まず概念の定義を一文で書き、次に講義の例へ対応付け、最後に意思決定モデルにどのような影響を与えるかを述べる。単語の列挙だけでは不十分であり、データがどのように情報モデルへ変換され、その情報モデルがどのように決定変数・目的関数・制約へ接続されるかまで説明する。
\item \textbf{Slide 017「準確率的モデリング」を、配送またはTSPの具体例で説明せよ。}\\旅行時間分布が分かるならstochastic、10-25分という範囲だけ分かるならquasi-stochasticとして扱う。 答案では、例を出した後に、それがどのモデル要素に対応するかを明示する。例えば、顧客集合や旅行時間行列は情報モデル、訪問順序や車両割当は意思決定モデル、実際の配送指示は実装である。
\end{enumerate}
\end{minipage}
\end{minipage}
\end{adjustbox}

\clearpage
\SlideHeader{Lecture 3: Uncertainty}{018}{データからシナリオを構築 I}
\SlideImage{18}{../Materials/2026/3DM_03_Uncertainty_SS26.pdf}
\begin{adjustbox}{max totalsize={\textwidth}{0.675\textheight},center}
\begin{minipage}{\textwidth}\scriptsize
\begin{minipage}[t]{0.485\textwidth}
\NoteSection{日本語訳}
\begin{itemize}
\item 旅行時間の単一観測の例。
\item Braunschweig, Hildesheimer Strasse, 月曜16-17時、n=20。
\end{itemize}
\NoteSection{試験対策ポイント}
\begin{itemize}
\item deterministic, stochastic, quasi-stochastic の違いは、将来値と確率情報をどう扱うかで決まる。
\item 確率分布がある場合とない場合で、使える評価基準が変わる。
\end{itemize}
\end{minipage}\hfill
\begin{minipage}[t]{0.485\textwidth}
\NoteSection{解説}
このページでは「データからシナリオを構築 I」を、講義全体の流れの中で位置づける。スライドの箇条書きは短いが、試験ではその背後にある概念、現実例、モデル上の意味を文章で説明できる必要がある。\par
Deterministic model は、入力値が既知で固定されていると仮定する。例えば、AからBへの旅行時間を常に15分と置く。これは単純で計算しやすいが、実際の交通変動やサービス時間のばらつきを無視するため、現実で計画が崩れることがある。\par
Stochastic model は、不確実な変数を確率分布として表す。旅行時間が平均20分・標準偏差5分の分布に従う、注文発生がポアソン過程に従う、サービス時間が対数正規分布に従う、といった表現である。確率が分かるため、期待費用、分散、分位点、確率制約、サンプリングを使える。\par
Quasi-stochastic model は、確率分布は分からないが、起こり得るシナリオや値の範囲は分かる場合に使う。例えば、旅行時間が10分から25分の間であることは分かるが、その確率は分からない場合である。このとき、区間、代表シナリオ、最悪ケース、Hurwicz基準、minmax regret が使われる。\par
確率分布が十分に推定できるなら stochastic model が自然である。しかし、過去データが少ない、新しいサービスで履歴がない、環境が変化して過去データが当てにならない場合は quasi-stochastic model が有用である。どちらを使うかは、データ量、分布仮定の妥当性、意思決定の目的によって決まる。\par
試験では、stochastic と quasi-stochastic の違いを『不確実であるかどうか』ではなく『確率情報を持っているかどうか』で説明する。どちらも不確実性を扱うが、stochastic は確率を使い、quasi-stochastic は確率なしで区間やシナリオを使う。\par
\NoteSection{確認問題と解答}
\begin{enumerate}
\item \textbf{Slide 018「データからシナリオを構築 I」の中心概念を、専門用語を使わずに説明せよ。}\\このスライドでは、データからシナリオを構築 I を通じて deterministic, stochastic, quasi-stochastic の違いは、将来値と確率情報をどう扱うかで決まる。 ことを理解する必要がある。試験答案では、まず概念の定義を一文で書き、次に講義の例へ対応付け、最後に意思決定モデルにどのような影響を与えるかを述べる。単語の列挙だけでは不十分であり、データがどのように情報モデルへ変換され、その情報モデルがどのように決定変数・目的関数・制約へ接続されるかまで説明する。
\item \textbf{Slide 018「データからシナリオを構築 I」を、配送またはTSPの具体例で説明せよ。}\\旅行時間分布が分かるならstochastic、10-25分という範囲だけ分かるならquasi-stochasticとして扱う。 答案では、例を出した後に、それがどのモデル要素に対応するかを明示する。例えば、顧客集合や旅行時間行列は情報モデル、訪問順序や車両割当は意思決定モデル、実際の配送指示は実装である。
\end{enumerate}
\end{minipage}
\end{minipage}
\end{adjustbox}

\clearpage
\SlideHeader{Lecture 3: Uncertainty}{019}{データからシナリオを構築 II}
\SlideImage{19}{../Materials/2026/3DM_03_Uncertainty_SS26.pdf}
\begin{adjustbox}{max totalsize={\textwidth}{0.675\textheight},center}
\begin{minipage}{\textwidth}\scriptsize
\begin{minipage}[t]{0.485\textwidth}
\NoteSection{日本語訳}
\begin{itemize}
\item 構築するシナリオ数を決める。例: 5。
\item 極端シナリオを選ぶ。最良ケース・最悪ケース。
\item 中間シナリオを選ぶ。
\item 選択はドメイン知識に基づく。
\item 一般にはシナリオ数を増やすと推定精度は高まる。
\item ただし意思決定の複雑性も増す。
\end{itemize}
\NoteSection{試験対策ポイント}
\begin{itemize}
\item deterministic, stochastic, quasi-stochastic の違いは、将来値と確率情報をどう扱うかで決まる。
\item 確率分布がある場合とない場合で、使える評価基準が変わる。
\end{itemize}
\end{minipage}\hfill
\begin{minipage}[t]{0.485\textwidth}
\NoteSection{解説}
このページでは「データからシナリオを構築 II」を、講義全体の流れの中で位置づける。スライドの箇条書きは短いが、試験ではその背後にある概念、現実例、モデル上の意味を文章で説明できる必要がある。\par
Deterministic model は、入力値が既知で固定されていると仮定する。例えば、AからBへの旅行時間を常に15分と置く。これは単純で計算しやすいが、実際の交通変動やサービス時間のばらつきを無視するため、現実で計画が崩れることがある。\par
Stochastic model は、不確実な変数を確率分布として表す。旅行時間が平均20分・標準偏差5分の分布に従う、注文発生がポアソン過程に従う、サービス時間が対数正規分布に従う、といった表現である。確率が分かるため、期待費用、分散、分位点、確率制約、サンプリングを使える。\par
Quasi-stochastic model は、確率分布は分からないが、起こり得るシナリオや値の範囲は分かる場合に使う。例えば、旅行時間が10分から25分の間であることは分かるが、その確率は分からない場合である。このとき、区間、代表シナリオ、最悪ケース、Hurwicz基準、minmax regret が使われる。\par
確率分布が十分に推定できるなら stochastic model が自然である。しかし、過去データが少ない、新しいサービスで履歴がない、環境が変化して過去データが当てにならない場合は quasi-stochastic model が有用である。どちらを使うかは、データ量、分布仮定の妥当性、意思決定の目的によって決まる。\par
試験では、stochastic と quasi-stochastic の違いを『不確実であるかどうか』ではなく『確率情報を持っているかどうか』で説明する。どちらも不確実性を扱うが、stochastic は確率を使い、quasi-stochastic は確率なしで区間やシナリオを使う。\par
\NoteSection{確認問題と解答}
\begin{enumerate}
\item \textbf{Slide 019「データからシナリオを構築 II」の中心概念を、専門用語を使わずに説明せよ。}\\このスライドでは、データからシナリオを構築 II を通じて deterministic, stochastic, quasi-stochastic の違いは、将来値と確率情報をどう扱うかで決まる。 ことを理解する必要がある。試験答案では、まず概念の定義を一文で書き、次に講義の例へ対応付け、最後に意思決定モデルにどのような影響を与えるかを述べる。単語の列挙だけでは不十分であり、データがどのように情報モデルへ変換され、その情報モデルがどのように決定変数・目的関数・制約へ接続されるかまで説明する。
\item \textbf{Slide 019「データからシナリオを構築 II」を、配送またはTSPの具体例で説明せよ。}\\旅行時間分布が分かるならstochastic、10-25分という範囲だけ分かるならquasi-stochasticとして扱う。 答案では、例を出した後に、それがどのモデル要素に対応するかを明示する。例えば、顧客集合や旅行時間行列は情報モデル、訪問順序や車両割当は意思決定モデル、実際の配送指示は実装である。
\end{enumerate}
\end{minipage}
\end{minipage}
\end{adjustbox}

\clearpage
\SlideHeader{Lecture 3: Uncertainty}{020}{データから区間を決める}
\SlideImage{20}{../Materials/2026/3DM_03_Uncertainty_SS26.pdf}
\begin{adjustbox}{max totalsize={\textwidth}{0.675\textheight},center}
\begin{minipage}{\textwidth}\scriptsize
\begin{minipage}[t]{0.485\textwidth}
\NoteSection{日本語訳}
\begin{itemize}
\item 限界
\begin{itemize}
\item 上限: 90\%分位点。
\item 下限: 10\%分位点。
\end{itemize}
\item 限界の調整
\begin{itemize}
\item 悲観的: 最悪ケース側へずらす。
\item 楽観的: 最良ケース側へずらす。
\item 区間幅を狭める、または広げる。
\end{itemize}
\end{itemize}
\NoteSection{試験対策ポイント}
\begin{itemize}
\item deterministic, stochastic, quasi-stochastic の違いは、将来値と確率情報をどう扱うかで決まる。
\item 確率分布がある場合とない場合で、使える評価基準が変わる。
\end{itemize}
\end{minipage}\hfill
\begin{minipage}[t]{0.485\textwidth}
\NoteSection{解説}
このページでは「データから区間を決める」を、講義全体の流れの中で位置づける。スライドの箇条書きは短いが、試験ではその背後にある概念、現実例、モデル上の意味を文章で説明できる必要がある。\par
Deterministic model は、入力値が既知で固定されていると仮定する。例えば、AからBへの旅行時間を常に15分と置く。これは単純で計算しやすいが、実際の交通変動やサービス時間のばらつきを無視するため、現実で計画が崩れることがある。\par
Stochastic model は、不確実な変数を確率分布として表す。旅行時間が平均20分・標準偏差5分の分布に従う、注文発生がポアソン過程に従う、サービス時間が対数正規分布に従う、といった表現である。確率が分かるため、期待費用、分散、分位点、確率制約、サンプリングを使える。\par
Quasi-stochastic model は、確率分布は分からないが、起こり得るシナリオや値の範囲は分かる場合に使う。例えば、旅行時間が10分から25分の間であることは分かるが、その確率は分からない場合である。このとき、区間、代表シナリオ、最悪ケース、Hurwicz基準、minmax regret が使われる。\par
確率分布が十分に推定できるなら stochastic model が自然である。しかし、過去データが少ない、新しいサービスで履歴がない、環境が変化して過去データが当てにならない場合は quasi-stochastic model が有用である。どちらを使うかは、データ量、分布仮定の妥当性、意思決定の目的によって決まる。\par
試験では、stochastic と quasi-stochastic の違いを『不確実であるかどうか』ではなく『確率情報を持っているかどうか』で説明する。どちらも不確実性を扱うが、stochastic は確率を使い、quasi-stochastic は確率なしで区間やシナリオを使う。\par
\NoteSection{確認問題と解答}
\begin{enumerate}
\item \textbf{Slide 020「データから区間を決める」の中心概念を、専門用語を使わずに説明せよ。}\\このスライドでは、データから区間を決める を通じて deterministic, stochastic, quasi-stochastic の違いは、将来値と確率情報をどう扱うかで決まる。 ことを理解する必要がある。試験答案では、まず概念の定義を一文で書き、次に講義の例へ対応付け、最後に意思決定モデルにどのような影響を与えるかを述べる。単語の列挙だけでは不十分であり、データがどのように情報モデルへ変換され、その情報モデルがどのように決定変数・目的関数・制約へ接続されるかまで説明する。
\item \textbf{Slide 020「データから区間を決める」を、配送またはTSPの具体例で説明せよ。}\\旅行時間分布が分かるならstochastic、10-25分という範囲だけ分かるならquasi-stochasticとして扱う。 答案では、例を出した後に、それがどのモデル要素に対応するかを明示する。例えば、顧客集合や旅行時間行列は情報モデル、訪問順序や車両割当は意思決定モデル、実際の配送指示は実装である。
\end{enumerate}
\end{minipage}
\end{minipage}
\end{adjustbox}

\clearpage
\ExamHeader{不確実性の種類とモデリング}
\ExamQuestion{SS24 Aufgabe 2a [3 Punkte]}
\ExamSubsection{問題内容}
Same-Day配送における不確実性を、需要、リソース、環境の三カテゴリについて一例ずつ挙げる問題。\par
\ExamSubsection{満点答案}
需要の不確実性として、注文がいつ発生するか、どの地域から発生するか、注文量がどれくらいかがある。リソースの不確実性として、ドライバーの利用可能性、車両故障、バッテリー残量、積載容量、ドライバーのスキルがある。環境の不確実性として、交通渋滞、駐車可能性、天候、道路工事、事故、通信ネットワーク障害がある。\par
満点を取るには、単に三語を列挙するのではなく、それぞれが配送計画へどう影響するかを説明する。需要は訪問先と負荷を変え、リソースは実行可能な割当を変え、環境は旅行時間や遅延リスクを変える。\par
\ExamSubsection{具体例による解説}
夕方に予想外の大量注文が入ると需要不確実性、登録ドライバーがアプリを切って稼働しないとリソース不確実性、突然の道路閉鎖で到着時間が変わると環境不確実性である。\par
\ExamSubsection{採点で落とさない要素}
\begin{itemize}
\item 三カテゴリをすべて挙げる
\item 各カテゴリに具体例
\item 計画への影響を説明
\end{itemize}
\vspace{2mm}\hrule\vspace{2mm}
\ExamQuestion{WS24/25 Aufgabe 3 [12 Punkte]}
\ExamSubsection{問題内容}
stochastic と quasi-stochastic の違い、および各手法の考え方を説明する問題。\par
\ExamSubsection{満点答案}
stochastic model では、不確実な実現値について確率分布またはシナリオ確率が既知または推定可能である。したがって、期待値、分散、確率制約、Monte Carlo sampling などを用いて意思決定を評価できる。例として、旅行時間がガンマ分布に従うと仮定し、その期待値や分位点を使ってルートを評価する場合がある。\par
quasi-stochastic model では、確率分布は分からないが、起こり得るシナリオや区間は分かる。例えば、旅行時間が [l,u] の範囲にあることだけが分かる場合である。このとき、Hurwicz基準、最悪ケース最適化、minmax regret、What-if分析を使う。\par
両者の違いは、確率を使って平均的・分布的に評価できるか、確率なしで代表シナリオや区間に対して頑健に評価するかである。データが十分なら stochastic、データが少ない・分布仮定が難しいなら quasi-stochastic が適する。\par
\ExamSubsection{具体例による解説}
旅行時間の過去観測が大量にあり、分布をフィットできるなら stochastic。新しい道路工事で過去データがなく、専門家が『10分から25分』という範囲だけを与えるなら quasi-stochastic。\par
\ExamSubsection{採点で落とさない要素}
\begin{itemize}
\item 確率既知/未知の違い
\item stochastic手法
\item quasi-stochastic手法
\item 適用条件
\item 具体例
\end{itemize}
\vspace{2mm}\hrule\vspace{2mm}

\clearpage
\SlideHeader{Lecture 3: Uncertainty}{021}{アジェンダ}
\SlideImage{21}{../Materials/2026/3DM_03_Uncertainty_SS26.pdf}
\begin{adjustbox}{max totalsize={\textwidth}{0.675\textheight},center}
\begin{minipage}{\textwidth}\scriptsize
\begin{minipage}[t]{0.485\textwidth}
\NoteSection{日本語訳}
\begin{itemize}
\item 不確実性。
\item 不確実性のモデリング。
\item 不確実性下の意思決定。
\item ケーススタディ。
\end{itemize}
\end{minipage}\hfill
\begin{minipage}[t]{0.485\textwidth}
\NoteSection{注記}
{\small このページは表紙・目次・連絡先などの事務情報であるため、解説と確認問題は付けない。}
\end{minipage}
\end{minipage}
\end{adjustbox}

\clearpage
\SlideHeader{Lecture 3: Uncertainty}{022}{例: ナップサック問題 I}
\SlideImage{22}{../Materials/2026/3DM_03_Uncertainty_SS26.pdf}
\begin{adjustbox}{max totalsize={\textwidth}{0.675\textheight},center}
\begin{minipage}{\textwidth}\scriptsize
\begin{minipage}[t]{0.485\textwidth}
\NoteSection{日本語訳}
\begin{itemize}
\item 泥棒がナップサックを持って家に入り、後で売るために物を盗む。
\item 各物品には価値と重さがある。
\item 物品の組合せの総重量はナップサック容量に収まらなければならない。
\item 売却価格の合計が最大となる組合せを探す。
\item 決定論的環境
\begin{itemize}
\item 可能な物品の正確な価値 v\_i と重さ w\_i を知っている。
\item ナップサック容量 W を正確に知っている。
\end{itemize}
\end{itemize}
\NoteSection{試験対策ポイント}
\begin{itemize}
\item 決定 x の後にランダム情報 xi が実現するため、期待値やリスク基準で評価する。
\item Monte Carlo sampling は期待値を標本平均で近似する方法である。
\end{itemize}
\end{minipage}\hfill
\begin{minipage}[t]{0.485\textwidth}
\NoteSection{解説}
このページでは「例: ナップサック問題 I」を、講義全体の流れの中で位置づける。スライドの箇条書きは短いが、試験ではその背後にある概念、現実例、モデル上の意味を文章で説明できる必要がある。\par
確率的意思決定では、意思決定者は現在の情報に基づいて決定 x を選ぶが、その後にランダム情報 xi が実現する。例えばルートを選んだ後に、実際の交通、サービス時間、注文キャンセルが判明する。したがって、費用 F(x,xi) は意思決定時点では確定していない。\par
このため、単一の実現値に対する費用を最小化するのではなく、期待費用 E[F(x,xi)]、分散、分位点、確率制約、リスク尺度などを使って評価する。リスク中立なら期待値中心でよいが、遅延に厳しいサービスでは分散や上位分位点も重要になる。\par
Monte Carlo Sampling は、確率分布から多数のシナリオを生成し、それぞれで費用を計算して平均する方法である。解析的に期待値を計算できない場合でも、サンプル平均によって近似できる。サンプル数が増えるほど推定は安定するが、計算量も増える。\par
期待値-分散原理では、平均性能とリスクを同時に見る。例えばルートAは平均45分だがばらつきが大きく、ルートBは平均50分だが安定している場合、時間厳守が重要ならBを選ぶ理由がある。平均だけではサービス信頼性を評価できない。\par
試験では、F(x,xi)を直接最適化できない理由を明確に述べる。xi は決定時点で未知のランダム変数であり、事前に分かるのは分布やサンプルだけである。そのため、期待値やリスクを使う、という流れで説明する。\par
\NoteSection{確認問題と解答}
\begin{enumerate}
\item \textbf{Slide 022「例: ナップサック問題 I」の中心概念を、専門用語を使わずに説明せよ。}\\このスライドでは、例: ナップサック問題 I を通じて 決定 x の後にランダム情報 xi が実現するため、期待値やリスク基準で評価する。 ことを理解する必要がある。試験答案では、まず概念の定義を一文で書き、次に講義の例へ対応付け、最後に意思決定モデルにどのような影響を与えるかを述べる。単語の列挙だけでは不十分であり、データがどのように情報モデルへ変換され、その情報モデルがどのように決定変数・目的関数・制約へ接続されるかまで説明する。
\item \textbf{Slide 022「例: ナップサック問題 I」を、配送またはTSPの具体例で説明せよ。}\\ルート選択後に実際の交通が実現するため、期待費用やサンプリング平均でルートを評価する。 答案では、例を出した後に、それがどのモデル要素に対応するかを明示する。例えば、顧客集合や旅行時間行列は情報モデル、訪問順序や車両割当は意思決定モデル、実際の配送指示は実装である。
\end{enumerate}
\end{minipage}
\end{minipage}
\end{adjustbox}

\clearpage
\SlideHeader{Lecture 3: Uncertainty}{023}{例: ナップサック問題 II}
\SlideImage{23}{../Materials/2026/3DM_03_Uncertainty_SS26.pdf}
\begin{adjustbox}{max totalsize={\textwidth}{0.675\textheight},center}
\begin{minipage}{\textwidth}\scriptsize
\begin{minipage}[t]{0.485\textwidth}
\NoteSection{日本語訳}
\begin{itemize}
\item 確率的環境
\begin{itemize}
\item 目的関数の不確実性: 闇市場での正確な売却価格が分からない。
\item 制約の不確実性: 物品の正確な重さが分からない。
\item 制約の不確実性: ナップサックの正確な容量が分からない。
\end{itemize}
\item 例
\begin{itemize}
\item 目的関数にランダム変数がある。
\item 制約左辺にランダム変数がある。
\item 制約右辺にランダム変数がある。
\end{itemize}
\end{itemize}
\NoteSection{試験対策ポイント}
\begin{itemize}
\item 決定 x の後にランダム情報 xi が実現するため、期待値やリスク基準で評価する。
\item Monte Carlo sampling は期待値を標本平均で近似する方法である。
\end{itemize}
\end{minipage}\hfill
\begin{minipage}[t]{0.485\textwidth}
\NoteSection{解説}
このページでは「例: ナップサック問題 II」を、講義全体の流れの中で位置づける。スライドの箇条書きは短いが、試験ではその背後にある概念、現実例、モデル上の意味を文章で説明できる必要がある。\par
確率的意思決定では、意思決定者は現在の情報に基づいて決定 x を選ぶが、その後にランダム情報 xi が実現する。例えばルートを選んだ後に、実際の交通、サービス時間、注文キャンセルが判明する。したがって、費用 F(x,xi) は意思決定時点では確定していない。\par
このため、単一の実現値に対する費用を最小化するのではなく、期待費用 E[F(x,xi)]、分散、分位点、確率制約、リスク尺度などを使って評価する。リスク中立なら期待値中心でよいが、遅延に厳しいサービスでは分散や上位分位点も重要になる。\par
Monte Carlo Sampling は、確率分布から多数のシナリオを生成し、それぞれで費用を計算して平均する方法である。解析的に期待値を計算できない場合でも、サンプル平均によって近似できる。サンプル数が増えるほど推定は安定するが、計算量も増える。\par
期待値-分散原理では、平均性能とリスクを同時に見る。例えばルートAは平均45分だがばらつきが大きく、ルートBは平均50分だが安定している場合、時間厳守が重要ならBを選ぶ理由がある。平均だけではサービス信頼性を評価できない。\par
試験では、F(x,xi)を直接最適化できない理由を明確に述べる。xi は決定時点で未知のランダム変数であり、事前に分かるのは分布やサンプルだけである。そのため、期待値やリスクを使う、という流れで説明する。\par
\NoteSection{確認問題と解答}
\begin{enumerate}
\item \textbf{Slide 023「例: ナップサック問題 II」の中心概念を、専門用語を使わずに説明せよ。}\\このスライドでは、例: ナップサック問題 II を通じて 決定 x の後にランダム情報 xi が実現するため、期待値やリスク基準で評価する。 ことを理解する必要がある。試験答案では、まず概念の定義を一文で書き、次に講義の例へ対応付け、最後に意思決定モデルにどのような影響を与えるかを述べる。単語の列挙だけでは不十分であり、データがどのように情報モデルへ変換され、その情報モデルがどのように決定変数・目的関数・制約へ接続されるかまで説明する。
\item \textbf{Slide 023「例: ナップサック問題 II」を、配送またはTSPの具体例で説明せよ。}\\ルート選択後に実際の交通が実現するため、期待費用やサンプリング平均でルートを評価する。 答案では、例を出した後に、それがどのモデル要素に対応するかを明示する。例えば、顧客集合や旅行時間行列は情報モデル、訪問順序や車両割当は意思決定モデル、実際の配送指示は実装である。
\end{enumerate}
\end{minipage}
\end{minipage}
\end{adjustbox}

\clearpage
\SlideHeader{Lecture 3: Uncertainty}{024}{不確実性下の意思決定}
\SlideImage{24}{../Materials/2026/3DM_03_Uncertainty_SS26.pdf}
\begin{adjustbox}{max totalsize={\textwidth}{0.675\textheight},center}
\begin{minipage}{\textwidth}\scriptsize
\begin{minipage}[t]{0.485\textwidth}
\NoteSection{日本語訳}
\begin{itemize}
\item 確率的意思決定は準確率的意思決定と異なる。
\item 不確実性は目的関数と制約のどちらにも影響し得る。
\item 実現値は決定後に現れる。
\item 確率的意思決定: 期待値と偏差が既知である。
\item 準確率的意思決定: 期待値と偏差が未知である。
\end{itemize}
\NoteSection{試験対策ポイント}
\begin{itemize}
\item deterministic, stochastic, quasi-stochastic の違いは、将来値と確率情報をどう扱うかで決まる。
\item 確率分布がある場合とない場合で、使える評価基準が変わる。
\end{itemize}
\end{minipage}\hfill
\begin{minipage}[t]{0.485\textwidth}
\NoteSection{解説}
このページでは「不確実性下の意思決定」を、講義全体の流れの中で位置づける。スライドの箇条書きは短いが、試験ではその背後にある概念、現実例、モデル上の意味を文章で説明できる必要がある。\par
Deterministic model は、入力値が既知で固定されていると仮定する。例えば、AからBへの旅行時間を常に15分と置く。これは単純で計算しやすいが、実際の交通変動やサービス時間のばらつきを無視するため、現実で計画が崩れることがある。\par
Stochastic model は、不確実な変数を確率分布として表す。旅行時間が平均20分・標準偏差5分の分布に従う、注文発生がポアソン過程に従う、サービス時間が対数正規分布に従う、といった表現である。確率が分かるため、期待費用、分散、分位点、確率制約、サンプリングを使える。\par
Quasi-stochastic model は、確率分布は分からないが、起こり得るシナリオや値の範囲は分かる場合に使う。例えば、旅行時間が10分から25分の間であることは分かるが、その確率は分からない場合である。このとき、区間、代表シナリオ、最悪ケース、Hurwicz基準、minmax regret が使われる。\par
確率分布が十分に推定できるなら stochastic model が自然である。しかし、過去データが少ない、新しいサービスで履歴がない、環境が変化して過去データが当てにならない場合は quasi-stochastic model が有用である。どちらを使うかは、データ量、分布仮定の妥当性、意思決定の目的によって決まる。\par
試験では、stochastic と quasi-stochastic の違いを『不確実であるかどうか』ではなく『確率情報を持っているかどうか』で説明する。どちらも不確実性を扱うが、stochastic は確率を使い、quasi-stochastic は確率なしで区間やシナリオを使う。\par
\NoteSection{確認問題と解答}
\begin{enumerate}
\item \textbf{Slide 024「不確実性下の意思決定」の中心概念を、専門用語を使わずに説明せよ。}\\このスライドでは、不確実性下の意思決定 を通じて deterministic, stochastic, quasi-stochastic の違いは、将来値と確率情報をどう扱うかで決まる。 ことを理解する必要がある。試験答案では、まず概念の定義を一文で書き、次に講義の例へ対応付け、最後に意思決定モデルにどのような影響を与えるかを述べる。単語の列挙だけでは不十分であり、データがどのように情報モデルへ変換され、その情報モデルがどのように決定変数・目的関数・制約へ接続されるかまで説明する。
\item \textbf{Slide 024「不確実性下の意思決定」を、配送またはTSPの具体例で説明せよ。}\\旅行時間分布が分かるならstochastic、10-25分という範囲だけ分かるならquasi-stochasticとして扱う。 答案では、例を出した後に、それがどのモデル要素に対応するかを明示する。例えば、顧客集合や旅行時間行列は情報モデル、訪問順序や車両割当は意思決定モデル、実際の配送指示は実装である。
\end{enumerate}
\end{minipage}
\end{minipage}
\end{adjustbox}

\clearpage
\SlideHeader{Lecture 3: Uncertainty}{025}{確率的意思決定}
\SlideImage{25}{../Materials/2026/3DM_03_Uncertainty_SS26.pdf}
\begin{adjustbox}{max totalsize={\textwidth}{0.675\textheight},center}
\begin{minipage}{\textwidth}\scriptsize
\begin{minipage}[t]{0.485\textwidth}
\NoteSection{日本語訳}
\begin{itemize}
\item X はすべての実行可能決定の領域である。
\item x は特定の決定である。
\item xi はランダム情報であり、決定後に利用可能になる。
\item F は最小化される費用関数であり、F(x,xi) と表される。
\item F(x,xi) を直接最適化することはできない。
\item 期待値 E[F(x,xi)] を最小化する。
\item 最適値と最適決定集合を定義する。
\end{itemize}
\NoteSection{試験対策ポイント}
\begin{itemize}
\item 決定 x の後にランダム情報 xi が実現するため、期待値やリスク基準で評価する。
\item Monte Carlo sampling は期待値を標本平均で近似する方法である。
\end{itemize}
\end{minipage}\hfill
\begin{minipage}[t]{0.485\textwidth}
\NoteSection{解説}
このページでは「確率的意思決定」を、講義全体の流れの中で位置づける。スライドの箇条書きは短いが、試験ではその背後にある概念、現実例、モデル上の意味を文章で説明できる必要がある。\par
確率的意思決定では、意思決定者は現在の情報に基づいて決定 x を選ぶが、その後にランダム情報 xi が実現する。例えばルートを選んだ後に、実際の交通、サービス時間、注文キャンセルが判明する。したがって、費用 F(x,xi) は意思決定時点では確定していない。\par
このため、単一の実現値に対する費用を最小化するのではなく、期待費用 E[F(x,xi)]、分散、分位点、確率制約、リスク尺度などを使って評価する。リスク中立なら期待値中心でよいが、遅延に厳しいサービスでは分散や上位分位点も重要になる。\par
Monte Carlo Sampling は、確率分布から多数のシナリオを生成し、それぞれで費用を計算して平均する方法である。解析的に期待値を計算できない場合でも、サンプル平均によって近似できる。サンプル数が増えるほど推定は安定するが、計算量も増える。\par
期待値-分散原理では、平均性能とリスクを同時に見る。例えばルートAは平均45分だがばらつきが大きく、ルートBは平均50分だが安定している場合、時間厳守が重要ならBを選ぶ理由がある。平均だけではサービス信頼性を評価できない。\par
試験では、F(x,xi)を直接最適化できない理由を明確に述べる。xi は決定時点で未知のランダム変数であり、事前に分かるのは分布やサンプルだけである。そのため、期待値やリスクを使う、という流れで説明する。\par
\NoteSection{確認問題と解答}
\begin{enumerate}
\item \textbf{Slide 025「確率的意思決定」の中心概念を、専門用語を使わずに説明せよ。}\\このスライドでは、確率的意思決定 を通じて 決定 x の後にランダム情報 xi が実現するため、期待値やリスク基準で評価する。 ことを理解する必要がある。試験答案では、まず概念の定義を一文で書き、次に講義の例へ対応付け、最後に意思決定モデルにどのような影響を与えるかを述べる。単語の列挙だけでは不十分であり、データがどのように情報モデルへ変換され、その情報モデルがどのように決定変数・目的関数・制約へ接続されるかまで説明する。
\item \textbf{Slide 025「確率的意思決定」を、配送またはTSPの具体例で説明せよ。}\\ルート選択後に実際の交通が実現するため、期待費用やサンプリング平均でルートを評価する。 答案では、例を出した後に、それがどのモデル要素に対応するかを明示する。例えば、顧客集合や旅行時間行列は情報モデル、訪問順序や車両割当は意思決定モデル、実際の配送指示は実装である。
\end{enumerate}
\end{minipage}
\end{minipage}
\end{adjustbox}

\clearpage
\SlideHeader{Lecture 3: Uncertainty}{026}{確率的意思決定 II: Monte Carlo Sampling}
\SlideImage{26}{../Materials/2026/3DM_03_Uncertainty_SS26.pdf}
\begin{adjustbox}{max totalsize={\textwidth}{0.675\textheight},center}
\begin{minipage}{\textwidth}\scriptsize
\begin{minipage}[t]{0.485\textwidth}
\NoteSection{日本語訳}
\begin{itemize}
\item ランダム情報 xi から n 個のシナリオ xi\_i を選ぶ。
\item 目的関数 E[F(x,xi)] を標本平均で近似する。
\item 右辺の決定論的問題を解く。
\item 近似最適値を求める。
\item 例
\begin{itemize}
\item 各物品の売却価格分布が既知である。
\item 各物品価格をサンプリングして n 個のシナリオを生成する。
\item 各シナリオを独立に解く。
\item n個の解から近似最適値と近似最適決定集合を決定する。
\end{itemize}
\end{itemize}
\NoteSection{試験対策ポイント}
\begin{itemize}
\item 決定 x の後にランダム情報 xi が実現するため、期待値やリスク基準で評価する。
\item Monte Carlo sampling は期待値を標本平均で近似する方法である。
\end{itemize}
\end{minipage}\hfill
\begin{minipage}[t]{0.485\textwidth}
\NoteSection{解説}
このページでは「確率的意思決定 II: Monte Carlo Sampling」を、講義全体の流れの中で位置づける。スライドの箇条書きは短いが、試験ではその背後にある概念、現実例、モデル上の意味を文章で説明できる必要がある。\par
確率的意思決定では、意思決定者は現在の情報に基づいて決定 x を選ぶが、その後にランダム情報 xi が実現する。例えばルートを選んだ後に、実際の交通、サービス時間、注文キャンセルが判明する。したがって、費用 F(x,xi) は意思決定時点では確定していない。\par
このため、単一の実現値に対する費用を最小化するのではなく、期待費用 E[F(x,xi)]、分散、分位点、確率制約、リスク尺度などを使って評価する。リスク中立なら期待値中心でよいが、遅延に厳しいサービスでは分散や上位分位点も重要になる。\par
Monte Carlo Sampling は、確率分布から多数のシナリオを生成し、それぞれで費用を計算して平均する方法である。解析的に期待値を計算できない場合でも、サンプル平均によって近似できる。サンプル数が増えるほど推定は安定するが、計算量も増える。\par
期待値-分散原理では、平均性能とリスクを同時に見る。例えばルートAは平均45分だがばらつきが大きく、ルートBは平均50分だが安定している場合、時間厳守が重要ならBを選ぶ理由がある。平均だけではサービス信頼性を評価できない。\par
試験では、F(x,xi)を直接最適化できない理由を明確に述べる。xi は決定時点で未知のランダム変数であり、事前に分かるのは分布やサンプルだけである。そのため、期待値やリスクを使う、という流れで説明する。\par
\NoteSection{確認問題と解答}
\begin{enumerate}
\item \textbf{Slide 026「確率的意思決定 II: Monte Carlo Sampling」の中心概念を、専門用語を使わずに説明せよ。}\\このスライドでは、確率的意思決定 II: Monte Carlo Sampling を通じて 決定 x の後にランダム情報 xi が実現するため、期待値やリスク基準で評価する。 ことを理解する必要がある。試験答案では、まず概念の定義を一文で書き、次に講義の例へ対応付け、最後に意思決定モデルにどのような影響を与えるかを述べる。単語の列挙だけでは不十分であり、データがどのように情報モデルへ変換され、その情報モデルがどのように決定変数・目的関数・制約へ接続されるかまで説明する。
\item \textbf{Slide 026「確率的意思決定 II: Monte Carlo Sampling」を、配送またはTSPの具体例で説明せよ。}\\ルート選択後に実際の交通が実現するため、期待費用やサンプリング平均でルートを評価する。 答案では、例を出した後に、それがどのモデル要素に対応するかを明示する。例えば、顧客集合や旅行時間行列は情報モデル、訪問順序や車両割当は意思決定モデル、実際の配送指示は実装である。
\end{enumerate}
\end{minipage}
\end{minipage}
\end{adjustbox}

\clearpage
\SlideHeader{Lecture 3: Uncertainty}{027}{確率的意思決定 III: 期待値-分散原理}
\SlideImage{27}{../Materials/2026/3DM_03_Uncertainty_SS26.pdf}
\begin{adjustbox}{max totalsize={\textwidth}{0.675\textheight},center}
\begin{minipage}{\textwidth}\scriptsize
\begin{minipage}[t]{0.485\textwidth}
\NoteSection{日本語訳}
\begin{itemize}
\item x は期待値と分散を用いる選好関数で評価される。
\item より精密な目的関数のために分散を組み込む。
\item 分散は q で重み付けされる: mu(x) + q sigma\textasciicircum{}2(x)。
\item q=-1: リスクへの親和性が高い。
\item q=0: リスク中立。
\item q=1: リスク回避。
\item 例
\begin{itemize}
\item 泥棒は各物品の売却価格分布を知っている。
\item 履歴データに基づく仮定など。
\item 期待値と分散を計算する。
\item 総期待売却価格が最もよい組合せを選ぶ。
\end{itemize}
\end{itemize}
\NoteSection{試験対策ポイント}
\begin{itemize}
\item 決定 x の後にランダム情報 xi が実現するため、期待値やリスク基準で評価する。
\item Monte Carlo sampling は期待値を標本平均で近似する方法である。
\end{itemize}
\end{minipage}\hfill
\begin{minipage}[t]{0.485\textwidth}
\NoteSection{解説}
このページでは「確率的意思決定 III: 期待値-分散原理」を、講義全体の流れの中で位置づける。スライドの箇条書きは短いが、試験ではその背後にある概念、現実例、モデル上の意味を文章で説明できる必要がある。\par
確率的意思決定では、意思決定者は現在の情報に基づいて決定 x を選ぶが、その後にランダム情報 xi が実現する。例えばルートを選んだ後に、実際の交通、サービス時間、注文キャンセルが判明する。したがって、費用 F(x,xi) は意思決定時点では確定していない。\par
このため、単一の実現値に対する費用を最小化するのではなく、期待費用 E[F(x,xi)]、分散、分位点、確率制約、リスク尺度などを使って評価する。リスク中立なら期待値中心でよいが、遅延に厳しいサービスでは分散や上位分位点も重要になる。\par
Monte Carlo Sampling は、確率分布から多数のシナリオを生成し、それぞれで費用を計算して平均する方法である。解析的に期待値を計算できない場合でも、サンプル平均によって近似できる。サンプル数が増えるほど推定は安定するが、計算量も増える。\par
期待値-分散原理では、平均性能とリスクを同時に見る。例えばルートAは平均45分だがばらつきが大きく、ルートBは平均50分だが安定している場合、時間厳守が重要ならBを選ぶ理由がある。平均だけではサービス信頼性を評価できない。\par
試験では、F(x,xi)を直接最適化できない理由を明確に述べる。xi は決定時点で未知のランダム変数であり、事前に分かるのは分布やサンプルだけである。そのため、期待値やリスクを使う、という流れで説明する。\par
\NoteSection{確認問題と解答}
\begin{enumerate}
\item \textbf{Slide 027「確率的意思決定 III: 期待値-分散原理」の中心概念を、専門用語を使わずに説明せよ。}\\このスライドでは、確率的意思決定 III: 期待値-分散原理 を通じて 決定 x の後にランダム情報 xi が実現するため、期待値やリスク基準で評価する。 ことを理解する必要がある。試験答案では、まず概念の定義を一文で書き、次に講義の例へ対応付け、最後に意思決定モデルにどのような影響を与えるかを述べる。単語の列挙だけでは不十分であり、データがどのように情報モデルへ変換され、その情報モデルがどのように決定変数・目的関数・制約へ接続されるかまで説明する。
\item \textbf{Slide 027「確率的意思決定 III: 期待値-分散原理」を、配送またはTSPの具体例で説明せよ。}\\ルート選択後に実際の交通が実現するため、期待費用やサンプリング平均でルートを評価する。 答案では、例を出した後に、それがどのモデル要素に対応するかを明示する。例えば、顧客集合や旅行時間行列は情報モデル、訪問順序や車両割当は意思決定モデル、実際の配送指示は実装である。
\end{enumerate}
\end{minipage}
\end{minipage}
\end{adjustbox}

\clearpage
\SlideHeader{Lecture 3: Uncertainty}{028}{準確率的意思決定 I: Hurwicz基準}
\SlideImage{28}{../Materials/2026/3DM_03_Uncertainty_SS26.pdf}
\begin{adjustbox}{max totalsize={\textwidth}{0.675\textheight},center}
\begin{minipage}{\textwidth}\scriptsize
\begin{minipage}[t]{0.485\textwidth}
\NoteSection{日本語訳}
\begin{itemize}
\item 最悪実現と最良実現を考慮する。
\item 楽観パラメータ lambda in [0,1] を用いて重み付き和を計算する。
\item (1-lambda) × worst case + lambda × best case。
\item 例
\begin{itemize}
\item 泥棒は各物品の過去の売却価格を知っている。
\item 異なるシナリオを選ぶ。
\item 各物品について重み付き平均または合計価格を計算する。
\item 決定された総価格が最良となる組合せを選ぶ。
\end{itemize}
\end{itemize}
\NoteSection{試験対策ポイント}
\begin{itemize}
\item Hurwicz は最良ケースと最悪ケースの重み付け、minmax regret は最大の後悔を最小化する。
\item regret は、実現シナリオを事前に知っていた場合の最適解との差である。
\end{itemize}
\end{minipage}\hfill
\begin{minipage}[t]{0.485\textwidth}
\NoteSection{解説}
このページでは「準確率的意思決定 I: Hurwicz基準」を、講義全体の流れの中で位置づける。スライドの箇条書きは短いが、試験ではその背後にある概念、現実例、モデル上の意味を文章で説明できる必要がある。\par
Quasi-stochastic な状況では、確率分布は分からないが、複数のシナリオや区間は分かる。例えば旅行時間が10-25分の範囲にあることは分かるが、10分になる確率、25分になる確率は分からない。このような場合、期待値を正確には計算できない。\par
Hurwicz基準は、最良ケースと最悪ケースを楽観係数で重み付けして評価する。楽観的な意思決定者は最良ケースを重視し、悲観的な意思決定者は最悪ケースを重視する。これは簡単だが、係数の選び方に主観が入る。\par
Minmax regret は、各シナリオで『そのシナリオを事前に知っていれば選べた最適解』との差を regret として測る。候補解ごとに全シナリオの regret を計算し、その最大 regret が最小の解を選ぶ。これは、最悪の後悔を抑える基準である。\par
regret の直感は、後から結果が分かったときに『別の解なら大幅に良かった』という損失を小さくすることである。単純な最悪ケース最適化よりも、各シナリオでの相対的な失敗を見ている点が特徴である。\par
具体例として、平均的に短いが渋滞時に極端に悪化するルートと、少し長いが安定したルートがある。minmax regret は、各交通シナリオで最良ルートとの差を比べ、最悪の差が小さい安定ルートを選ぶ可能性がある。\par
\NoteSection{確認問題と解答}
\begin{enumerate}
\item \textbf{Slide 028「準確率的意思決定 I: Hurwicz基準」の中心概念を、専門用語を使わずに説明せよ。}\\このスライドでは、準確率的意思決定 I: Hurwicz基準 を通じて Hurwicz は最良ケースと最悪ケースの重み付け、minmax regret は最大の後悔を最小化する。 ことを理解する必要がある。試験答案では、まず概念の定義を一文で書き、次に講義の例へ対応付け、最後に意思決定モデルにどのような影響を与えるかを述べる。単語の列挙だけでは不十分であり、データがどのように情報モデルへ変換され、その情報モデルがどのように決定変数・目的関数・制約へ接続されるかまで説明する。
\item \textbf{Slide 028「準確率的意思決定 I: Hurwicz基準」を、配送またはTSPの具体例で説明せよ。}\\各交通シナリオで最良ルートとの差をregretとして計算し、最大regretが小さいルートを選ぶ。 答案では、例を出した後に、それがどのモデル要素に対応するかを明示する。例えば、顧客集合や旅行時間行列は情報モデル、訪問順序や車両割当は意思決定モデル、実際の配送指示は実装である。
\end{enumerate}
\end{minipage}
\end{minipage}
\end{adjustbox}

\clearpage
\SlideHeader{Lecture 3: Uncertainty}{029}{準確率的意思決定 II: Minmax-Regret基準}
\SlideImage{29}{../Materials/2026/3DM_03_Uncertainty_SS26.pdf}
\begin{adjustbox}{max totalsize={\textwidth}{0.675\textheight},center}
\begin{minipage}{\textwidth}\scriptsize
\begin{minipage}[t]{0.485\textwidth}
\NoteSection{日本語訳}
\begin{itemize}
\item あるシナリオにおける悲観的決定を、同じシナリオで完全情報がある場合の最良決定と比較する。
\item Regret は、同じシナリオにおける最良決定からの乖離である。
\item 例
\begin{itemize}
\item 泥棒が最良ケースを実現できなかったことへのregret。
\item 泥棒は各物品の過去の売却価格を知っている。
\item 各組合せについて最良ケースと最悪ケースの差を計算する。
\item 取った物品については最悪売却価格、残した物品については最良売却価格を考える。
\item 安定した市場価格の物品を取り、楽観ケースでだけ高価な物品を残す結果になり得る。
\end{itemize}
\end{itemize}
\NoteSection{試験対策ポイント}
\begin{itemize}
\item Hurwicz は最良ケースと最悪ケースの重み付け、minmax regret は最大の後悔を最小化する。
\item regret は、実現シナリオを事前に知っていた場合の最適解との差である。
\end{itemize}
\end{minipage}\hfill
\begin{minipage}[t]{0.485\textwidth}
\NoteSection{解説}
このページでは「準確率的意思決定 II: Minmax-Regret基準」を、講義全体の流れの中で位置づける。スライドの箇条書きは短いが、試験ではその背後にある概念、現実例、モデル上の意味を文章で説明できる必要がある。\par
Quasi-stochastic な状況では、確率分布は分からないが、複数のシナリオや区間は分かる。例えば旅行時間が10-25分の範囲にあることは分かるが、10分になる確率、25分になる確率は分からない。このような場合、期待値を正確には計算できない。\par
Hurwicz基準は、最良ケースと最悪ケースを楽観係数で重み付けして評価する。楽観的な意思決定者は最良ケースを重視し、悲観的な意思決定者は最悪ケースを重視する。これは簡単だが、係数の選び方に主観が入る。\par
Minmax regret は、各シナリオで『そのシナリオを事前に知っていれば選べた最適解』との差を regret として測る。候補解ごとに全シナリオの regret を計算し、その最大 regret が最小の解を選ぶ。これは、最悪の後悔を抑える基準である。\par
regret の直感は、後から結果が分かったときに『別の解なら大幅に良かった』という損失を小さくすることである。単純な最悪ケース最適化よりも、各シナリオでの相対的な失敗を見ている点が特徴である。\par
具体例として、平均的に短いが渋滞時に極端に悪化するルートと、少し長いが安定したルートがある。minmax regret は、各交通シナリオで最良ルートとの差を比べ、最悪の差が小さい安定ルートを選ぶ可能性がある。\par
\NoteSection{確認問題と解答}
\begin{enumerate}
\item \textbf{Slide 029「準確率的意思決定 II: Minmax-Regret基準」の中心概念を、専門用語を使わずに説明せよ。}\\このスライドでは、準確率的意思決定 II: Minmax-Regret基準 を通じて Hurwicz は最良ケースと最悪ケースの重み付け、minmax regret は最大の後悔を最小化する。 ことを理解する必要がある。試験答案では、まず概念の定義を一文で書き、次に講義の例へ対応付け、最後に意思決定モデルにどのような影響を与えるかを述べる。単語の列挙だけでは不十分であり、データがどのように情報モデルへ変換され、その情報モデルがどのように決定変数・目的関数・制約へ接続されるかまで説明する。
\item \textbf{Slide 029「準確率的意思決定 II: Minmax-Regret基準」を、配送またはTSPの具体例で説明せよ。}\\各交通シナリオで最良ルートとの差をregretとして計算し、最大regretが小さいルートを選ぶ。 答案では、例を出した後に、それがどのモデル要素に対応するかを明示する。例えば、顧客集合や旅行時間行列は情報モデル、訪問順序や車両割当は意思決定モデル、実際の配送指示は実装である。
\end{enumerate}
\end{minipage}
\end{minipage}
\end{adjustbox}

\clearpage
\ExamHeader{不確実性下の意思決定: expectation, Hurwicz, regret}
\ExamQuestion{SS24 Aufgabe 2b [3 Punkte]}
\ExamSubsection{問題内容}
修理サービスや在宅介護の顧客サービス時間について、確率密度関数の典型的性質を三つ説明する問題。\par
\ExamSubsection{満点答案}
第一に、サービス時間は負にならないため、分布の定義域は0以上である。第二に、多くのサービスは標準的な短い所要時間に集中するが、例外的に長い作業が起こるため右裾が長くなりやすい。第三に、サービス品質や遅延リスクを評価するには平均だけでは不十分であり、分散、標準偏差、分位点、特に上位分位点が重要である。\par
試験では、単に『正規分布』など分布名を書くよりも、なぜ非負・右裾・ばらつきが意思決定に影響するかを説明する方がよい。サービス時間の上振れは次顧客への遅延を引き起こし、ツアー全体に伝播する。\par
\ExamSubsection{具体例による解説}
家電修理では多くの訪問は20分で終わるが、故障が複雑な場合は90分かかる。この長い右裾を無視して平均30分だけで計画すると、後続訪問が連鎖的に遅れる。\par
\ExamSubsection{採点で落とさない要素}
\begin{itemize}
\item 非負性
\item 右裾/歪み
\item 平均だけでなく分散・分位点
\item 計画への影響
\end{itemize}
\vspace{2mm}\hrule\vspace{2mm}
\ExamQuestion{SS23 Aufgabe 3 [6 Punkte] / WS24/25 Aufgabe 1 [10 Punkte]}
\ExamSubsection{問題内容}
区間旅行時間、楽観値・悲観値、regret に基づき、最大の後悔を抑える意思決定の論理を説明する問題。\par
\ExamSubsection{満点答案}
区間旅行時間では、各アークについて楽観的旅行時間 l\_\{ij\} と悲観的旅行時間 u\_\{ij\} が分かるが、その確率は分からない。この場合、平均値で最適化するのではなく、複数の可能な交通状態に対して解を評価する。regret とは、あるシナリオが実現したときに、自分が選んだ解の費用と、そのシナリオを事前に知っていた場合の最適解の費用との差である。\par
minmax regret では、各候補ルートについて全シナリオの regret を計算し、その最大値を求める。そして最大 regret が最小になるルートを選ぶ。これにより、どのシナリオが実現しても『別のルートなら大幅に良かった』という最悪の後悔を小さくできる。\par
この基準は、最良ケースだけを見てリスクを無視することも、最悪ケースだけを見て過度に保守的になることも避ける。都市物流では、効率性とサービス信頼性のバランスを取るために有用である。\par
\ExamSubsection{具体例による解説}
ルートAは通常なら短いが、渋滞時に大きく悪化する。ルートBは通常少し長いが、どの状態でも安定している。minmax regretでは、シナリオごとに完全情報最適ルートとの差を見て、最悪の差が小さいBを選ぶ可能性がある。\par
\ExamSubsection{採点で落とさない要素}
\begin{itemize}
\item regretの定義
\item 完全情報最適解との差
\item 最大regretを最小化
\item 平均/最悪ケースとの違い
\item 物流例
\end{itemize}
\vspace{2mm}\hrule\vspace{2mm}

\clearpage
\SlideHeader{Lecture 3: Uncertainty}{030}{アジェンダ}
\SlideImage{30}{../Materials/2026/3DM_03_Uncertainty_SS26.pdf}
\begin{adjustbox}{max totalsize={\textwidth}{0.675\textheight},center}
\begin{minipage}{\textwidth}\scriptsize
\begin{minipage}[t]{0.485\textwidth}
\NoteSection{日本語訳}
\begin{itemize}
\item 不確実性。
\item 不確実性のモデリング。
\item 不確実性下の意思決定。
\item ケーススタディ。
\end{itemize}
\end{minipage}\hfill
\begin{minipage}[t]{0.485\textwidth}
\NoteSection{注記}
{\small このページは表紙・目次・連絡先などの事務情報であるため、解説と確認問題は付けない。}
\end{minipage}
\end{minipage}
\end{adjustbox}

\clearpage
\SlideHeader{Lecture 3: Uncertainty}{031}{ケーススタディ: 都市物流のロバスト車両ルーティング}
\SlideImage{31}{../Materials/2026/3DM_03_Uncertainty_SS26.pdf}
\begin{adjustbox}{max totalsize={\textwidth}{0.675\textheight},center}
\begin{minipage}{\textwidth}\scriptsize
\begin{minipage}[t]{0.485\textwidth}
\NoteSection{日本語訳}
\begin{itemize}
\item 都市物流の課題
\begin{itemize}
\item 高い需要とコスト圧力。
\item 高い顧客期待。
\end{itemize}
\item 都市交通システム
\begin{itemize}
\item 限定されたインフラ: 渋滞。
\item 変動しやすい旅行時間: 不確実性。
\end{itemize}
\item 都市物流の車両ルーティングへ、旅行時間不確実性の情報を統合する。
\end{itemize}
\NoteSection{試験対策ポイント}
\begin{itemize}
\item ロバスト解は、単一平均ケースで最短ではなく、複数シナリオで大きく崩れにくい解である。
\item 都市物流では効率性と信頼性のトレードオフが重要である。
\end{itemize}
\end{minipage}\hfill
\begin{minipage}[t]{0.485\textwidth}
\NoteSection{解説}
このページでは「ケーススタディ: 都市物流のロバスト車両ルーティング」を、講義全体の流れの中で位置づける。スライドの箇条書きは短いが、試験ではその背後にある概念、現実例、モデル上の意味を文章で説明できる必要がある。\par
都市物流では、旅行時間の変動、狭い時間窓、駐車問題、顧客の待ち時間、車両数制約が重なるため、平均旅行時間だけを最小化する計画では不十分なことが多い。平均では良くても、渋滞や遅いサービスが起きると多くの顧客に遅延が伝播する。\par
区間旅行時間 [l,u] は、各アークの旅行時間が下限 l と上限 u の間にあると表す方法である。確率分布を正確に推定できない場合でも、過去データの分位点や専門知識から現実的な範囲を作ることができる。これは準確率的情報モデルの典型である。\par
ロバストなルートは、平均ケースで必ず最短とは限らない。しかし、複数の交通シナリオや悪条件に対して大きく悪化しにくい。顧客サービスや時間約束が重要な場合、少し長いが安定したルートの方が、実務上は優れていることがある。\par
一方で、ロバスト性を重視しすぎると過度に保守的な解になる。すべての最悪ケースだけを想定すると、余裕を取りすぎてコストが大きくなる可能性がある。したがって、効率性、信頼性、後悔、サービスレベルのバランスを考える必要がある。\par
試験では、ロバスト解を『どんな場合でも絶対に最適な解』とは書かない。正しくは、複数の不確実なシナリオに対して性能が大きく崩れにくい解である。平均最適、最悪ケース最適、minmax regret の違いも説明できるとよい。\par
\NoteSection{確認問題と解答}
\begin{enumerate}
\item \textbf{Slide 031「ケーススタディ: 都市物流のロバスト車両ルーティング」の中心概念を、専門用語を使わずに説明せよ。}\\このスライドでは、ケーススタディ: 都市物流のロバスト車両ルーティング を通じて ロバスト解は、単一平均ケースで最短ではなく、複数シナリオで大きく崩れにくい解である。 ことを理解する必要がある。試験答案では、まず概念の定義を一文で書き、次に講義の例へ対応付け、最後に意思決定モデルにどのような影響を与えるかを述べる。単語の列挙だけでは不十分であり、データがどのように情報モデルへ変換され、その情報モデルがどのように決定変数・目的関数・制約へ接続されるかまで説明する。
\item \textbf{Slide 031「ケーススタディ: 都市物流のロバスト車両ルーティング」を、配送またはTSPの具体例で説明せよ。}\\平均では少し長くても、渋滞時に大きく遅れないルートはロバストなルートといえる。 答案では、例を出した後に、それがどのモデル要素に対応するかを明示する。例えば、顧客集合や旅行時間行列は情報モデル、訪問順序や車両割当は意思決定モデル、実際の配送指示は実装である。
\end{enumerate}
\end{minipage}
\end{minipage}
\end{adjustbox}

\clearpage
\SlideHeader{Lecture 3: Uncertainty}{032}{都市物流における役割}
\SlideImage{32}{../Materials/2026/3DM_03_Uncertainty_SS26.pdf}
\begin{adjustbox}{max totalsize={\textwidth}{0.675\textheight},center}
\begin{minipage}{\textwidth}\scriptsize
\begin{minipage}[t]{0.485\textwidth}
\NoteSection{日本語訳}
\begin{itemize}
\item Shipping Company
\begin{itemize}
\item オンライン小売業者。
\item 既知の顧客数。
\item 商品の配送を物流サービスプロバイダに委託する。
\item 約束された配送日に基づき alpha サービスレベルを要求する。
\end{itemize}
\item Logistic service provider
\begin{itemize}
\item 小包サービス。
\item 配送車両フリートを運用する。
\item 事前または前日計画を行う。
\item 旅行時間を最小化する。
\end{itemize}
\item alpha サービスレベルは旅行時間の影響を受ける。
\item 旅行時間の不確実性に関する情報は計画段階で考慮されなければならない。
\end{itemize}
\NoteSection{試験対策ポイント}
\begin{itemize}
\item ロバスト解は、単一平均ケースで最短ではなく、複数シナリオで大きく崩れにくい解である。
\item 都市物流では効率性と信頼性のトレードオフが重要である。
\end{itemize}
\end{minipage}\hfill
\begin{minipage}[t]{0.485\textwidth}
\NoteSection{解説}
このページでは「都市物流における役割」を、講義全体の流れの中で位置づける。スライドの箇条書きは短いが、試験ではその背後にある概念、現実例、モデル上の意味を文章で説明できる必要がある。\par
都市物流では、旅行時間の変動、狭い時間窓、駐車問題、顧客の待ち時間、車両数制約が重なるため、平均旅行時間だけを最小化する計画では不十分なことが多い。平均では良くても、渋滞や遅いサービスが起きると多くの顧客に遅延が伝播する。\par
区間旅行時間 [l,u] は、各アークの旅行時間が下限 l と上限 u の間にあると表す方法である。確率分布を正確に推定できない場合でも、過去データの分位点や専門知識から現実的な範囲を作ることができる。これは準確率的情報モデルの典型である。\par
ロバストなルートは、平均ケースで必ず最短とは限らない。しかし、複数の交通シナリオや悪条件に対して大きく悪化しにくい。顧客サービスや時間約束が重要な場合、少し長いが安定したルートの方が、実務上は優れていることがある。\par
一方で、ロバスト性を重視しすぎると過度に保守的な解になる。すべての最悪ケースだけを想定すると、余裕を取りすぎてコストが大きくなる可能性がある。したがって、効率性、信頼性、後悔、サービスレベルのバランスを考える必要がある。\par
試験では、ロバスト解を『どんな場合でも絶対に最適な解』とは書かない。正しくは、複数の不確実なシナリオに対して性能が大きく崩れにくい解である。平均最適、最悪ケース最適、minmax regret の違いも説明できるとよい。\par
\NoteSection{確認問題と解答}
\begin{enumerate}
\item \textbf{Slide 032「都市物流における役割」の中心概念を、専門用語を使わずに説明せよ。}\\このスライドでは、都市物流における役割 を通じて ロバスト解は、単一平均ケースで最短ではなく、複数シナリオで大きく崩れにくい解である。 ことを理解する必要がある。試験答案では、まず概念の定義を一文で書き、次に講義の例へ対応付け、最後に意思決定モデルにどのような影響を与えるかを述べる。単語の列挙だけでは不十分であり、データがどのように情報モデルへ変換され、その情報モデルがどのように決定変数・目的関数・制約へ接続されるかまで説明する。
\item \textbf{Slide 032「都市物流における役割」を、配送またはTSPの具体例で説明せよ。}\\平均では少し長くても、渋滞時に大きく遅れないルートはロバストなルートといえる。 答案では、例を出した後に、それがどのモデル要素に対応するかを明示する。例えば、顧客集合や旅行時間行列は情報モデル、訪問順序や車両割当は意思決定モデル、実際の配送指示は実装である。
\end{enumerate}
\end{minipage}
\end{minipage}
\end{adjustbox}

\clearpage
\SlideHeader{Lecture 3: Uncertainty}{033}{旅行時間不確実性のモデル化}
\SlideImage{33}{../Materials/2026/3DM_03_Uncertainty_SS26.pdf}
\begin{adjustbox}{max totalsize={\textwidth}{0.675\textheight},center}
\begin{minipage}{\textwidth}\scriptsize
\begin{minipage}[t]{0.485\textwidth}
\NoteSection{日本語訳}
\begin{itemize}
\item Quasi-stochastic
\begin{itemize}
\item 上限・下限 [l,u] を持つ区間旅行時間。
\item 少数の観測で利用しやすい。
\item 区間限界を決定する。
\item 区間演算。
\item 直感的で単純。
\end{itemize}
\item Stochastic
\begin{itemize}
\item 期待値 mu と分散 sigma\textasciicircum{}2 を持つ旅行時間分布。
\item 多数の観測を必要とする。
\item 分布型とパラメータ化が必要。
\item 畳み込みによる集約。
\item 抽象的で複雑。
\end{itemize}
\item 都市物流には区間旅行時間が適切な情報モデルである。
\end{itemize}
\NoteSection{試験対策ポイント}
\begin{itemize}
\item ロバスト解は、単一平均ケースで最短ではなく、複数シナリオで大きく崩れにくい解である。
\item 都市物流では効率性と信頼性のトレードオフが重要である。
\end{itemize}
\end{minipage}\hfill
\begin{minipage}[t]{0.485\textwidth}
\NoteSection{解説}
このページでは「旅行時間不確実性のモデル化」を、講義全体の流れの中で位置づける。スライドの箇条書きは短いが、試験ではその背後にある概念、現実例、モデル上の意味を文章で説明できる必要がある。\par
都市物流では、旅行時間の変動、狭い時間窓、駐車問題、顧客の待ち時間、車両数制約が重なるため、平均旅行時間だけを最小化する計画では不十分なことが多い。平均では良くても、渋滞や遅いサービスが起きると多くの顧客に遅延が伝播する。\par
区間旅行時間 [l,u] は、各アークの旅行時間が下限 l と上限 u の間にあると表す方法である。確率分布を正確に推定できない場合でも、過去データの分位点や専門知識から現実的な範囲を作ることができる。これは準確率的情報モデルの典型である。\par
ロバストなルートは、平均ケースで必ず最短とは限らない。しかし、複数の交通シナリオや悪条件に対して大きく悪化しにくい。顧客サービスや時間約束が重要な場合、少し長いが安定したルートの方が、実務上は優れていることがある。\par
一方で、ロバスト性を重視しすぎると過度に保守的な解になる。すべての最悪ケースだけを想定すると、余裕を取りすぎてコストが大きくなる可能性がある。したがって、効率性、信頼性、後悔、サービスレベルのバランスを考える必要がある。\par
試験では、ロバスト解を『どんな場合でも絶対に最適な解』とは書かない。正しくは、複数の不確実なシナリオに対して性能が大きく崩れにくい解である。平均最適、最悪ケース最適、minmax regret の違いも説明できるとよい。\par
\NoteSection{確認問題と解答}
\begin{enumerate}
\item \textbf{Slide 033「旅行時間不確実性のモデル化」の中心概念を、専門用語を使わずに説明せよ。}\\このスライドでは、旅行時間不確実性のモデル化 を通じて ロバスト解は、単一平均ケースで最短ではなく、複数シナリオで大きく崩れにくい解である。 ことを理解する必要がある。試験答案では、まず概念の定義を一文で書き、次に講義の例へ対応付け、最後に意思決定モデルにどのような影響を与えるかを述べる。単語の列挙だけでは不十分であり、データがどのように情報モデルへ変換され、その情報モデルがどのように決定変数・目的関数・制約へ接続されるかまで説明する。
\item \textbf{Slide 033「旅行時間不確実性のモデル化」を、配送またはTSPの具体例で説明せよ。}\\平均では少し長くても、渋滞時に大きく遅れないルートはロバストなルートといえる。 答案では、例を出した後に、それがどのモデル要素に対応するかを明示する。例えば、顧客集合や旅行時間行列は情報モデル、訪問順序や車両割当は意思決定モデル、実際の配送指示は実装である。
\end{enumerate}
\end{minipage}
\end{minipage}
\end{adjustbox}

\clearpage
\SlideHeader{Lecture 3: Uncertainty}{034}{情報モデルデータの生成}
\SlideImage{34}{../Materials/2026/3DM_03_Uncertainty_SS26.pdf}
\begin{adjustbox}{max totalsize={\textwidth}{0.675\textheight},center}
\begin{minipage}{\textwidth}\scriptsize
\begin{minipage}[t]{0.485\textwidth}
\NoteSection{日本語訳}
\begin{itemize}
\item Shifted Gamma distribution。
\item パラメータ
\begin{itemize}
\item 平均旅行時間。
\item 変動係数。
\end{itemize}
\item サンプリングアプローチ。
\item 500個の観測をランダムに抽出する。
\item 接続の例: 平均旅行時間15.5、変動係数10\%と30\%。
\end{itemize}
\NoteSection{試験対策ポイント}
\begin{itemize}
\item 期待値は平均性能、分散や分位点はリスクを表す。
\item 分布フィットでは可視化、パラメータ推定、適合度確認が必要である。
\end{itemize}
\end{minipage}\hfill
\begin{minipage}[t]{0.485\textwidth}
\NoteSection{解説}
このページでは「情報モデルデータの生成」を、講義全体の流れの中で位置づける。スライドの箇条書きは短いが、試験ではその背後にある概念、現実例、モデル上の意味を文章で説明できる必要がある。\par
確率分布は、不確実な変数がどの値をどれくらいの確率で取るかを表す。配送では、旅行時間、サービス時間、注文間隔、需要量、キャンセル数などを分布で表せる。分布を使うことで、単一の平均値ではなく、ばらつきや極端な値も考慮できる。\par
期待値は長期的な平均を表すが、期待値だけでは不十分なことが多い。平均20分の旅行時間でも、ほとんどが18-22分なのか、10分と50分が混在しているのかでは、計画上のリスクが全く違う。分散や標準偏差は、このばらつきの大きさを表す。\par
分位点はサービスレベルを考える際に重要である。例えば90\%分位点が35分なら、90\%のケースで35分以内に到着することを意味する。顧客に到着時間を約束する場合、平均よりも上位分位点を見る方が安全な計画につながる。\par
実データから分布を決めるには、まず観測値を集め、外れ値や欠損を確認し、ヒストグラムや箱ひげ図で形を見る。その後、正規分布、対数正規分布、ガンマ分布など候補を選び、平均や分散などのパラメータを推定し、適合度を確認する。\par
サービス時間のように負にならない値では、正規分布よりも右裾を持つ非負分布が自然な場合がある。家電修理や在宅介護では、多くの作業は短時間で終わるが、例外的に長い作業が起こる。試験では、分布名だけでなく、なぜその形が現実に合うのかを説明する。\par
\NoteSection{確認問題と解答}
\begin{enumerate}
\item \textbf{Slide 034「情報モデルデータの生成」の中心概念を、専門用語を使わずに説明せよ。}\\このスライドでは、情報モデルデータの生成 を通じて 期待値は平均性能、分散や分位点はリスクを表す。 ことを理解する必要がある。試験答案では、まず概念の定義を一文で書き、次に講義の例へ対応付け、最後に意思決定モデルにどのような影響を与えるかを述べる。単語の列挙だけでは不十分であり、データがどのように情報モデルへ変換され、その情報モデルがどのように決定変数・目的関数・制約へ接続されるかまで説明する。
\item \textbf{Slide 034「情報モデルデータの生成」を、配送またはTSPの具体例で説明せよ。}\\旅行時間の平均だけでなく、標準偏差や90\%分位点を見ることで遅延リスクを評価する。 答案では、例を出した後に、それがどのモデル要素に対応するかを明示する。例えば、顧客集合や旅行時間行列は情報モデル、訪問順序や車両割当は意思決定モデル、実際の配送指示は実装である。
\end{enumerate}
\end{minipage}
\end{minipage}
\end{adjustbox}

\clearpage
\SlideHeader{Lecture 3: Uncertainty}{035}{情報モデル}
\SlideImage{35}{../Materials/2026/3DM_03_Uncertainty_SS26.pdf}
\begin{adjustbox}{max totalsize={\textwidth}{0.675\textheight},center}
\begin{minipage}{\textwidth}\scriptsize
\begin{minipage}[t]{0.485\textwidth}
\NoteSection{日本語訳}
\begin{itemize}
\item 二つの旅行時間行列。
\item U: 区間旅行時間の上限。ピーク交通の95\%分位点。
\item L: 区間旅行時間の下限。オフピーク交通の5\%分位点。
\item l\_ij <= u\_ij。
\end{itemize}
\NoteSection{試験対策ポイント}
\begin{itemize}
\item ロバスト解は、単一平均ケースで最短ではなく、複数シナリオで大きく崩れにくい解である。
\item 都市物流では効率性と信頼性のトレードオフが重要である。
\end{itemize}
\end{minipage}\hfill
\begin{minipage}[t]{0.485\textwidth}
\NoteSection{解説}
このページでは「情報モデル」を、講義全体の流れの中で位置づける。スライドの箇条書きは短いが、試験ではその背後にある概念、現実例、モデル上の意味を文章で説明できる必要がある。\par
都市物流では、旅行時間の変動、狭い時間窓、駐車問題、顧客の待ち時間、車両数制約が重なるため、平均旅行時間だけを最小化する計画では不十分なことが多い。平均では良くても、渋滞や遅いサービスが起きると多くの顧客に遅延が伝播する。\par
区間旅行時間 [l,u] は、各アークの旅行時間が下限 l と上限 u の間にあると表す方法である。確率分布を正確に推定できない場合でも、過去データの分位点や専門知識から現実的な範囲を作ることができる。これは準確率的情報モデルの典型である。\par
ロバストなルートは、平均ケースで必ず最短とは限らない。しかし、複数の交通シナリオや悪条件に対して大きく悪化しにくい。顧客サービスや時間約束が重要な場合、少し長いが安定したルートの方が、実務上は優れていることがある。\par
一方で、ロバスト性を重視しすぎると過度に保守的な解になる。すべての最悪ケースだけを想定すると、余裕を取りすぎてコストが大きくなる可能性がある。したがって、効率性、信頼性、後悔、サービスレベルのバランスを考える必要がある。\par
試験では、ロバスト解を『どんな場合でも絶対に最適な解』とは書かない。正しくは、複数の不確実なシナリオに対して性能が大きく崩れにくい解である。平均最適、最悪ケース最適、minmax regret の違いも説明できるとよい。\par
\NoteSection{確認問題と解答}
\begin{enumerate}
\item \textbf{Slide 035「情報モデル」の中心概念を、専門用語を使わずに説明せよ。}\\このスライドでは、情報モデル を通じて ロバスト解は、単一平均ケースで最短ではなく、複数シナリオで大きく崩れにくい解である。 ことを理解する必要がある。試験答案では、まず概念の定義を一文で書き、次に講義の例へ対応付け、最後に意思決定モデルにどのような影響を与えるかを述べる。単語の列挙だけでは不十分であり、データがどのように情報モデルへ変換され、その情報モデルがどのように決定変数・目的関数・制約へ接続されるかまで説明する。
\item \textbf{Slide 035「情報モデル」を、配送またはTSPの具体例で説明せよ。}\\平均では少し長くても、渋滞時に大きく遅れないルートはロバストなルートといえる。 答案では、例を出した後に、それがどのモデル要素に対応するかを明示する。例えば、顧客集合や旅行時間行列は情報モデル、訪問順序や車両割当は意思決定モデル、実際の配送指示は実装である。
\end{enumerate}
\end{minipage}
\end{minipage}
\end{adjustbox}

\clearpage
\SlideHeader{Lecture 3: Uncertainty}{036}{Minmax Regret基準を持つRobust TSP}
\SlideImage{36}{../Materials/2026/3DM_03_Uncertainty_SS26.pdf}
\begin{adjustbox}{max totalsize={\textwidth}{0.675\textheight},center}
\begin{minipage}{\textwidth}\scriptsize
\begin{minipage}[t]{0.485\textwidth}
\NoteSection{日本語訳}
\begin{itemize}
\item ツアーで使う辺。
\item 使わない辺。
\item Minmax regret により、複数の旅行時間実現に対して後悔の大きいルートを避ける。
\end{itemize}
\NoteSection{試験対策ポイント}
\begin{itemize}
\item Hurwicz は最良ケースと最悪ケースの重み付け、minmax regret は最大の後悔を最小化する。
\item regret は、実現シナリオを事前に知っていた場合の最適解との差である。
\end{itemize}
\end{minipage}\hfill
\begin{minipage}[t]{0.485\textwidth}
\NoteSection{解説}
このページでは「Minmax Regret基準を持つRobust TSP」を、講義全体の流れの中で位置づける。スライドの箇条書きは短いが、試験ではその背後にある概念、現実例、モデル上の意味を文章で説明できる必要がある。\par
Quasi-stochastic な状況では、確率分布は分からないが、複数のシナリオや区間は分かる。例えば旅行時間が10-25分の範囲にあることは分かるが、10分になる確率、25分になる確率は分からない。このような場合、期待値を正確には計算できない。\par
Hurwicz基準は、最良ケースと最悪ケースを楽観係数で重み付けして評価する。楽観的な意思決定者は最良ケースを重視し、悲観的な意思決定者は最悪ケースを重視する。これは簡単だが、係数の選び方に主観が入る。\par
Minmax regret は、各シナリオで『そのシナリオを事前に知っていれば選べた最適解』との差を regret として測る。候補解ごとに全シナリオの regret を計算し、その最大 regret が最小の解を選ぶ。これは、最悪の後悔を抑える基準である。\par
regret の直感は、後から結果が分かったときに『別の解なら大幅に良かった』という損失を小さくすることである。単純な最悪ケース最適化よりも、各シナリオでの相対的な失敗を見ている点が特徴である。\par
具体例として、平均的に短いが渋滞時に極端に悪化するルートと、少し長いが安定したルートがある。minmax regret は、各交通シナリオで最良ルートとの差を比べ、最悪の差が小さい安定ルートを選ぶ可能性がある。\par
\NoteSection{確認問題と解答}
\begin{enumerate}
\item \textbf{Slide 036「Minmax Regret基準を持つRobust TSP」の中心概念を、専門用語を使わずに説明せよ。}\\このスライドでは、Minmax Regret基準を持つRobust TSP を通じて Hurwicz は最良ケースと最悪ケースの重み付け、minmax regret は最大の後悔を最小化する。 ことを理解する必要がある。試験答案では、まず概念の定義を一文で書き、次に講義の例へ対応付け、最後に意思決定モデルにどのような影響を与えるかを述べる。単語の列挙だけでは不十分であり、データがどのように情報モデルへ変換され、その情報モデルがどのように決定変数・目的関数・制約へ接続されるかまで説明する。
\item \textbf{Slide 036「Minmax Regret基準を持つRobust TSP」を、配送またはTSPの具体例で説明せよ。}\\各交通シナリオで最良ルートとの差をregretとして計算し、最大regretが小さいルートを選ぶ。 答案では、例を出した後に、それがどのモデル要素に対応するかを明示する。例えば、顧客集合や旅行時間行列は情報モデル、訪問順序や車両割当は意思決定モデル、実際の配送指示は実装である。
\end{enumerate}
\end{minipage}
\end{minipage}
\end{adjustbox}

\clearpage
\SlideHeader{Lecture 3: Uncertainty}{037}{何を達成したいか}
\SlideImage{37}{../Materials/2026/3DM_03_Uncertainty_SS26.pdf}
\begin{adjustbox}{max totalsize={\textwidth}{0.675\textheight},center}
\begin{minipage}{\textwidth}\scriptsize
\begin{minipage}[t]{0.485\textwidth}
\NoteSection{日本語訳}
\begin{itemize}
\item 多くの交通速度シナリオで実装できるルートを見つけたい。
\item 旅行時間という効率性と、時間厳守というサービス信頼性のトレードオフを扱う。
\item 最良ケースだけを考えることは選択肢ではない。実行可能性を破る。
\item 道路セグメントごとの最良・最悪速度の平均はサービス品質を悪化させ得る。
\item 最悪ケース速度を考えると効率は悪いがサービス品質は良い。
\item regret に基づくアプローチが効率性と信頼性のバランスを取ることを期待する。
\item 評価: 自己設定した到着時刻からの早着・遅着の合計偏差。
\end{itemize}
\NoteSection{試験対策ポイント}
\begin{itemize}
\item ロバスト解は、単一平均ケースで最短ではなく、複数シナリオで大きく崩れにくい解である。
\item 都市物流では効率性と信頼性のトレードオフが重要である。
\end{itemize}
\end{minipage}\hfill
\begin{minipage}[t]{0.485\textwidth}
\NoteSection{解説}
このページでは「何を達成したいか」を、講義全体の流れの中で位置づける。スライドの箇条書きは短いが、試験ではその背後にある概念、現実例、モデル上の意味を文章で説明できる必要がある。\par
都市物流では、旅行時間の変動、狭い時間窓、駐車問題、顧客の待ち時間、車両数制約が重なるため、平均旅行時間だけを最小化する計画では不十分なことが多い。平均では良くても、渋滞や遅いサービスが起きると多くの顧客に遅延が伝播する。\par
区間旅行時間 [l,u] は、各アークの旅行時間が下限 l と上限 u の間にあると表す方法である。確率分布を正確に推定できない場合でも、過去データの分位点や専門知識から現実的な範囲を作ることができる。これは準確率的情報モデルの典型である。\par
ロバストなルートは、平均ケースで必ず最短とは限らない。しかし、複数の交通シナリオや悪条件に対して大きく悪化しにくい。顧客サービスや時間約束が重要な場合、少し長いが安定したルートの方が、実務上は優れていることがある。\par
一方で、ロバスト性を重視しすぎると過度に保守的な解になる。すべての最悪ケースだけを想定すると、余裕を取りすぎてコストが大きくなる可能性がある。したがって、効率性、信頼性、後悔、サービスレベルのバランスを考える必要がある。\par
試験では、ロバスト解を『どんな場合でも絶対に最適な解』とは書かない。正しくは、複数の不確実なシナリオに対して性能が大きく崩れにくい解である。平均最適、最悪ケース最適、minmax regret の違いも説明できるとよい。\par
\NoteSection{確認問題と解答}
\begin{enumerate}
\item \textbf{Slide 037「何を達成したいか」の中心概念を、専門用語を使わずに説明せよ。}\\このスライドでは、何を達成したいか を通じて ロバスト解は、単一平均ケースで最短ではなく、複数シナリオで大きく崩れにくい解である。 ことを理解する必要がある。試験答案では、まず概念の定義を一文で書き、次に講義の例へ対応付け、最後に意思決定モデルにどのような影響を与えるかを述べる。単語の列挙だけでは不十分であり、データがどのように情報モデルへ変換され、その情報モデルがどのように決定変数・目的関数・制約へ接続されるかまで説明する。
\item \textbf{Slide 037「何を達成したいか」を、配送またはTSPの具体例で説明せよ。}\\平均では少し長くても、渋滞時に大きく遅れないルートはロバストなルートといえる。 答案では、例を出した後に、それがどのモデル要素に対応するかを明示する。例えば、顧客集合や旅行時間行列は情報モデル、訪問順序や車両割当は意思決定モデル、実際の配送指示は実装である。
\end{enumerate}
\end{minipage}
\end{minipage}
\end{adjustbox}

\clearpage
\SlideHeader{Lecture 3: Uncertainty}{038}{問題設定}
\SlideImage{38}{../Materials/2026/3DM_03_Uncertainty_SS26.pdf}
\begin{adjustbox}{max totalsize={\textwidth}{0.675\textheight},center}
\begin{minipage}{\textwidth}\scriptsize
\begin{minipage}[t]{0.485\textwidth}
\NoteSection{日本語訳}
\begin{itemize}
\item 二階層配送ネットワーク。
\item サテライトは積み替え地点である。
\item デポから積み替え地点へ車両群で配送する。
\item 最終配送は各積み替え地点から始まる。
\item 目的: 積み替え地点での待機時間を含む旅行時間を最小化する。
\end{itemize}
\NoteSection{試験対策ポイント}
\begin{itemize}
\item ロバスト解は、単一平均ケースで最短ではなく、複数シナリオで大きく崩れにくい解である。
\item 都市物流では効率性と信頼性のトレードオフが重要である。
\end{itemize}
\end{minipage}\hfill
\begin{minipage}[t]{0.485\textwidth}
\NoteSection{解説}
このページでは「問題設定」を、講義全体の流れの中で位置づける。スライドの箇条書きは短いが、試験ではその背後にある概念、現実例、モデル上の意味を文章で説明できる必要がある。\par
都市物流では、旅行時間の変動、狭い時間窓、駐車問題、顧客の待ち時間、車両数制約が重なるため、平均旅行時間だけを最小化する計画では不十分なことが多い。平均では良くても、渋滞や遅いサービスが起きると多くの顧客に遅延が伝播する。\par
区間旅行時間 [l,u] は、各アークの旅行時間が下限 l と上限 u の間にあると表す方法である。確率分布を正確に推定できない場合でも、過去データの分位点や専門知識から現実的な範囲を作ることができる。これは準確率的情報モデルの典型である。\par
ロバストなルートは、平均ケースで必ず最短とは限らない。しかし、複数の交通シナリオや悪条件に対して大きく悪化しにくい。顧客サービスや時間約束が重要な場合、少し長いが安定したルートの方が、実務上は優れていることがある。\par
一方で、ロバスト性を重視しすぎると過度に保守的な解になる。すべての最悪ケースだけを想定すると、余裕を取りすぎてコストが大きくなる可能性がある。したがって、効率性、信頼性、後悔、サービスレベルのバランスを考える必要がある。\par
試験では、ロバスト解を『どんな場合でも絶対に最適な解』とは書かない。正しくは、複数の不確実なシナリオに対して性能が大きく崩れにくい解である。平均最適、最悪ケース最適、minmax regret の違いも説明できるとよい。\par
\NoteSection{確認問題と解答}
\begin{enumerate}
\item \textbf{Slide 038「問題設定」の中心概念を、専門用語を使わずに説明せよ。}\\このスライドでは、問題設定 を通じて ロバスト解は、単一平均ケースで最短ではなく、複数シナリオで大きく崩れにくい解である。 ことを理解する必要がある。試験答案では、まず概念の定義を一文で書き、次に講義の例へ対応付け、最後に意思決定モデルにどのような影響を与えるかを述べる。単語の列挙だけでは不十分であり、データがどのように情報モデルへ変換され、その情報モデルがどのように決定変数・目的関数・制約へ接続されるかまで説明する。
\item \textbf{Slide 038「問題設定」を、配送またはTSPの具体例で説明せよ。}\\平均では少し長くても、渋滞時に大きく遅れないルートはロバストなルートといえる。 答案では、例を出した後に、それがどのモデル要素に対応するかを明示する。例えば、顧客集合や旅行時間行列は情報モデル、訪問順序や車両割当は意思決定モデル、実際の配送指示は実装である。
\end{enumerate}
\end{minipage}
\end{minipage}
\end{adjustbox}

\clearpage
\SlideHeader{Lecture 3: Uncertainty}{039}{結果}
\SlideImage{39}{../Materials/2026/3DM_03_Uncertainty_SS26.pdf}
\begin{adjustbox}{max totalsize={\textwidth}{0.675\textheight},center}
\begin{minipage}{\textwidth}\scriptsize
\begin{minipage}[t]{0.485\textwidth}
\NoteSection{日本語訳}
\begin{itemize}
\item 手法
\begin{itemize}
\item minmax regret 基準を用いるVRP。
\item メタヒューリスティック。
\end{itemize}
\item 評価
\begin{itemize}
\item AVG: 平均。
\item UB: 最悪ケース。
\item ITT: 区間旅行時間。
\end{itemize}
\item 評価基準
\begin{itemize}
\item 使用車両数。
\item 総旅行時間。
\item サテライト到着の平均偏差、すなわち非同期化。
\end{itemize}
\end{itemize}
\NoteSection{試験対策ポイント}
\begin{itemize}
\item ロバスト解は、単一平均ケースで最短ではなく、複数シナリオで大きく崩れにくい解である。
\item 都市物流では効率性と信頼性のトレードオフが重要である。
\end{itemize}
\end{minipage}\hfill
\begin{minipage}[t]{0.485\textwidth}
\NoteSection{解説}
このページでは「結果」を、講義全体の流れの中で位置づける。スライドの箇条書きは短いが、試験ではその背後にある概念、現実例、モデル上の意味を文章で説明できる必要がある。\par
都市物流では、旅行時間の変動、狭い時間窓、駐車問題、顧客の待ち時間、車両数制約が重なるため、平均旅行時間だけを最小化する計画では不十分なことが多い。平均では良くても、渋滞や遅いサービスが起きると多くの顧客に遅延が伝播する。\par
区間旅行時間 [l,u] は、各アークの旅行時間が下限 l と上限 u の間にあると表す方法である。確率分布を正確に推定できない場合でも、過去データの分位点や専門知識から現実的な範囲を作ることができる。これは準確率的情報モデルの典型である。\par
ロバストなルートは、平均ケースで必ず最短とは限らない。しかし、複数の交通シナリオや悪条件に対して大きく悪化しにくい。顧客サービスや時間約束が重要な場合、少し長いが安定したルートの方が、実務上は優れていることがある。\par
一方で、ロバスト性を重視しすぎると過度に保守的な解になる。すべての最悪ケースだけを想定すると、余裕を取りすぎてコストが大きくなる可能性がある。したがって、効率性、信頼性、後悔、サービスレベルのバランスを考える必要がある。\par
試験では、ロバスト解を『どんな場合でも絶対に最適な解』とは書かない。正しくは、複数の不確実なシナリオに対して性能が大きく崩れにくい解である。平均最適、最悪ケース最適、minmax regret の違いも説明できるとよい。\par
\NoteSection{確認問題と解答}
\begin{enumerate}
\item \textbf{Slide 039「結果」の中心概念を、専門用語を使わずに説明せよ。}\\このスライドでは、結果 を通じて ロバスト解は、単一平均ケースで最短ではなく、複数シナリオで大きく崩れにくい解である。 ことを理解する必要がある。試験答案では、まず概念の定義を一文で書き、次に講義の例へ対応付け、最後に意思決定モデルにどのような影響を与えるかを述べる。単語の列挙だけでは不十分であり、データがどのように情報モデルへ変換され、その情報モデルがどのように決定変数・目的関数・制約へ接続されるかまで説明する。
\item \textbf{Slide 039「結果」を、配送またはTSPの具体例で説明せよ。}\\平均では少し長くても、渋滞時に大きく遅れないルートはロバストなルートといえる。 答案では、例を出した後に、それがどのモデル要素に対応するかを明示する。例えば、顧客集合や旅行時間行列は情報モデル、訪問順序や車両割当は意思決定モデル、実際の配送指示は実装である。
\end{enumerate}
\end{minipage}
\end{minipage}
\end{adjustbox}

\clearpage
\SlideHeader{Lecture 3: Uncertainty}{040}{まとめ}
\SlideImage{40}{../Materials/2026/3DM_03_Uncertainty_SS26.pdf}
\begin{adjustbox}{max totalsize={\textwidth}{0.675\textheight},center}
\begin{minipage}{\textwidth}\scriptsize
\begin{minipage}[t]{0.485\textwidth}
\NoteSection{日本語訳}
\begin{itemize}
\item 不確実性は情報モデルと意思決定モデルで考慮されるべきである。
\item 不確実性は多くの応用に影響し、意思決定に大きな影響を及ぼし得る。
\item 意思決定に関連するデータを有効に使うには、不確実性を適切にモデル化する必要がある。
\item 意思決定モデルは、応用と不確実性のモデリングに応じて選択されなければならない。
\end{itemize}
\NoteSection{試験対策ポイント}
\begin{itemize}
\item deterministic, stochastic, quasi-stochastic の違いは、将来値と確率情報をどう扱うかで決まる。
\item 確率分布がある場合とない場合で、使える評価基準が変わる。
\end{itemize}
\end{minipage}\hfill
\begin{minipage}[t]{0.485\textwidth}
\NoteSection{解説}
このページでは「まとめ」を、講義全体の流れの中で位置づける。スライドの箇条書きは短いが、試験ではその背後にある概念、現実例、モデル上の意味を文章で説明できる必要がある。\par
Deterministic model は、入力値が既知で固定されていると仮定する。例えば、AからBへの旅行時間を常に15分と置く。これは単純で計算しやすいが、実際の交通変動やサービス時間のばらつきを無視するため、現実で計画が崩れることがある。\par
Stochastic model は、不確実な変数を確率分布として表す。旅行時間が平均20分・標準偏差5分の分布に従う、注文発生がポアソン過程に従う、サービス時間が対数正規分布に従う、といった表現である。確率が分かるため、期待費用、分散、分位点、確率制約、サンプリングを使える。\par
Quasi-stochastic model は、確率分布は分からないが、起こり得るシナリオや値の範囲は分かる場合に使う。例えば、旅行時間が10分から25分の間であることは分かるが、その確率は分からない場合である。このとき、区間、代表シナリオ、最悪ケース、Hurwicz基準、minmax regret が使われる。\par
確率分布が十分に推定できるなら stochastic model が自然である。しかし、過去データが少ない、新しいサービスで履歴がない、環境が変化して過去データが当てにならない場合は quasi-stochastic model が有用である。どちらを使うかは、データ量、分布仮定の妥当性、意思決定の目的によって決まる。\par
試験では、stochastic と quasi-stochastic の違いを『不確実であるかどうか』ではなく『確率情報を持っているかどうか』で説明する。どちらも不確実性を扱うが、stochastic は確率を使い、quasi-stochastic は確率なしで区間やシナリオを使う。\par
\NoteSection{確認問題と解答}
\begin{enumerate}
\item \textbf{Slide 040「まとめ」の中心概念を、専門用語を使わずに説明せよ。}\\このスライドでは、まとめ を通じて deterministic, stochastic, quasi-stochastic の違いは、将来値と確率情報をどう扱うかで決まる。 ことを理解する必要がある。試験答案では、まず概念の定義を一文で書き、次に講義の例へ対応付け、最後に意思決定モデルにどのような影響を与えるかを述べる。単語の列挙だけでは不十分であり、データがどのように情報モデルへ変換され、その情報モデルがどのように決定変数・目的関数・制約へ接続されるかまで説明する。
\item \textbf{Slide 040「まとめ」を、配送またはTSPの具体例で説明せよ。}\\旅行時間分布が分かるならstochastic、10-25分という範囲だけ分かるならquasi-stochasticとして扱う。 答案では、例を出した後に、それがどのモデル要素に対応するかを明示する。例えば、顧客集合や旅行時間行列は情報モデル、訪問順序や車両割当は意思決定モデル、実際の配送指示は実装である。
\end{enumerate}
\end{minipage}
\end{minipage}
\end{adjustbox}

\clearpage
\ExamHeader{ロバスト都市物流と区間旅行時間}
\ExamQuestion{WS22/23 Aufgabe 5 [6 Punkte]}
\ExamSubsection{問題内容}
旅行時間不確実性を確率的または準確率的にモデル化し、ロバストな解が何を意味するかを説明する問題。\par
\ExamSubsection{満点答案}
確率的モデルでは、旅行時間の確率分布やシナリオ確率を仮定し、期待旅行時間、分散、分位点、サンプリングに基づいてルートを評価する。準確率的モデルでは、確率は分からないが、区間やシナリオ集合を使って不確実性を表す。区間旅行時間 [l\_\{ij\},u\_\{ij\}] はその典型である。\par
ロバストな解とは、単一の平均ケースで最も短いだけでなく、複数のあり得る交通状態で大きく性能が崩れにくい解である。例えば、平均旅行時間最小ルートはラッシュ時に大きく遅れる可能性がある。一方、ロバストルートは少し長くても、遅延リスクやサービス品質低下を抑える。\par
都市物流では、効率性と信頼性のトレードオフが重要である。最悪ケースだけに最適化すると過度に保守的になり、平均ケースだけではサービス違反が増える。regretや区間ベースの基準は、その中間を狙う。\par
\ExamSubsection{具体例による解説}
Aルートは平均45分だが渋滞時90分、Bルートは平均55分だが渋滞時65分なら、平均だけならA、ロバスト性ならBが有利になる場合がある。配送時間約束がある場合、Bの方が顧客サービスに適する。\par
\ExamSubsection{採点で落とさない要素}
\begin{itemize}
\item 確率的/準確率的の違い
\item 区間旅行時間
\item ロバスト解の定義
\item 効率性と信頼性のトレードオフ
\item 具体例
\end{itemize}
\vspace{2mm}\hrule\vspace{2mm}

\clearpage
\SlideHeader{Lecture 3: Uncertainty}{041}{さらなる問い}
\SlideImage{41}{../Materials/2026/3DM_03_Uncertainty_SS26.pdf}
\begin{adjustbox}{max totalsize={\textwidth}{0.675\textheight},center}
\begin{minipage}{\textwidth}\scriptsize
\begin{minipage}[t]{0.485\textwidth}
\NoteSection{日本語訳}
\begin{itemize}
\item 変化に反応できるか。
\item 例: 現在の交通情報に基づいてTSPを再計画する。これは dynamic decision making である。
\item 将来の不確実な発展を予測できるか。
\item 例: 渋滞リスクのある地域を避ける。
\item Anticipation。
\item Stochastic dynamic decision making。
\end{itemize}
\NoteSection{試験対策ポイント}
\begin{itemize}
\item 不確実性は、需要・リソース・環境の変化として整理できる。
\item 頻度、範囲、破壊力、予測可能性で性質が異なる。
\end{itemize}
\end{minipage}\hfill
\begin{minipage}[t]{0.485\textwidth}
\NoteSection{解説}
このページでは「さらなる問い」を、講義全体の流れの中で位置づける。スライドの箇条書きは短いが、試験ではその背後にある概念、現実例、モデル上の意味を文章で説明できる必要がある。\par
不確実性とは、意思決定時点では将来の状態や入力値が完全には分からないことである。配送であれば、注文がどこから来るか、交通がどれだけ混むか、サービス時間が何分か、ドライバーが利用可能かどうかが事前には確定していない。\par
不確実性は一種類ではない。需要不確実性は、顧客からの注文の発生時刻、場所、量に関する不確実性である。リソース不確実性は、車両故障、ドライバーのキャンセル、バッテリー残量、積載能力に関する不確実性である。環境不確実性は、交通、天候、道路工事、事故、駐車可能性に関する不確実性である。\par
不確実性は、頻度、範囲、破壊力、予測可能性によって性質が異なる。例えば日々の交通変動は頻繁に起きるが、過去データからある程度予測できる。一方で、大規模事故や道路閉鎖は低頻度だが、発生すれば計画を大きく壊す。\par
不確実性を無視して平均値だけで計画すると、平均的には良く見えるが実行時に頻繁に遅れる計画になる可能性がある。特に複数顧客を連続して訪問するルーティングでは、一つの遅延が後続訪問に波及するため、単独の平均時間以上にリスクが大きい。\par
試験答案では、不確実性の例を挙げるだけでなく、それが情報モデルと意思決定モデルにどう影響するかを書くとよい。情報モデルでは確率分布や区間として表現し、意思決定モデルでは期待値、リスク、ロバスト性、再最適化などを考える。\par
\NoteSection{確認問題と解答}
\begin{enumerate}
\item \textbf{Slide 041「さらなる問い」の中心概念を、専門用語を使わずに説明せよ。}\\このスライドでは、さらなる問い を通じて 不確実性は、需要・リソース・環境の変化として整理できる。 ことを理解する必要がある。試験答案では、まず概念の定義を一文で書き、次に講義の例へ対応付け、最後に意思決定モデルにどのような影響を与えるかを述べる。単語の列挙だけでは不十分であり、データがどのように情報モデルへ変換され、その情報モデルがどのように決定変数・目的関数・制約へ接続されるかまで説明する。
\item \textbf{Slide 041「さらなる問い」を、配送またはTSPの具体例で説明せよ。}\\注文がいつ来るか、ドライバーが稼働するか、道路が混むかは意思決定時点では確定しない。これが不確実性である。 答案では、例を出した後に、それがどのモデル要素に対応するかを明示する。例えば、顧客集合や旅行時間行列は情報モデル、訪問順序や車両割当は意思決定モデル、実際の配送指示は実装である。
\end{enumerate}
\end{minipage}
\end{minipage}
\end{adjustbox}
